\documentclass[11pt,a4paper]{article}
\usepackage[utf8]{inputenc}
\usepackage[T1]{fontenc}
\usepackage[french]{babel}
\usepackage{amsmath,amssymb,amsthm}
\usepackage{geometry}
\geometry{margin=1in}
\usepackage{tikz}
\usetikzlibrary{shapes.geometric, arrows, positioning}
\usepackage{listings}
\usepackage{xcolor}

\lstset{
    language=Python,
    basicstyle=\ttfamily\small,
    keywordstyle=\color{blue},
    stringstyle=\color{red},
    commentstyle=\color{green!50!black},
    showstringspaces=false,
    frame=single
}

\title{Logique formelle, connecteurs, tables de vérité et calcul des propositions}
\author{Charles EDOU NZE \\ \small Independent Researcher}
\date{\today}

\begin{document}
\maketitle

\begin{abstract}
Ce jalon introduit les fondements de la logique formelle, en explorant la syntaxe et la sémantique du calcul des propositions. À travers une approche rigoureuse et des démonstrations sans ellipse, il aborde les tables de vérité, les lois fondamentales (De Morgan, Peirce), et culmine avec une ouverture sur le SAT-solving en intelligence artificielle. Le jalon inclut également des exercices d'application théoriques et des travaux pratiques d'implémentation en Python.
\end{abstract}

\tableofcontents
\newpage

\part{Cours magistral et fondations théoriques}
\hypertarget{jalon-1-logique-formelle-connecteurs-tables-de-vuxe9rituxe9-et-calcul-des-propositions}{%
\section{Jalon 1 : Logique formelle, connecteurs, tables de vérité et
calcul des
propositions}\label{jalon-1-logique-formelle-connecteurs-tables-de-vuxe9rituxe9-et-calcul-des-propositions}}

\hypertarget{pruxe9sentation-du-concept-cluxe9}{%
\subsection{1. Présentation du concept
clé}\label{pruxe9sentation-du-concept-cluxe9}}

La logique formelle est la science des structures de raisonnement
valides. Pour aborder ce concept de manière géométrique et intuitive,
libérée du formalisme brut, imaginons un labyrinthe géant représentant
l'ensemble des possibles. Ce labyrinthe est composé d'une multitude
d'aiguillages et de portes commandés à distance. Chaque porte ne peut
être que dans deux états stables : - ouverte (valeur de vérité Vrai,
notée \(1\)) ou - fermée (valeur de vérité Faux, notée \(0\)).

La métaphore fondamentale de la logique propositionnelle réside dans la
notion de flot ou de flux d'information à travers un réseau
d'aiguillages. Considérons un fluide circulant dans ce réseau.

Si nous plaçons deux vannes à la suite sur une même canalisation, le
fluide ne passera que si la première vanne \textbf{ET} la seconde vanne
sont ouvertes.

C'est l'essence du connecteur de conjonction \textbf{(\(\land\))}. Si
nous plaçons ces deux vannes en parallèle sur deux embranchements
distincts reliant les mêmes points, il suffit que la première vanne OU
la seconde vanne soit ouverte pour que le fluide traverse. C'est
l'essence du connecteur de disjonction (\(\lor\)).


\begin{figure}[htbp]
\centering
\begin{tikzpicture}[
    node distance=2cm and 3cm,
    valve/.style={rectangle, draw=blue!60, fill=blue!5, very thick, minimum size=1cm},
    line/.style={draw, -latex', very thick}
]

% Nodes for AND logic (Series)
\node[valve] (v1) {Vanne $A$};
\node[valve, right=of v1] (v2) {Vanne $B$};
\node[left=1.5cm of v1] (in_and) {Flux entrant};
\node[right=1.5cm of v2] (out_and) {Flux sortant};

% Paths for AND
\path[line] (in_and) -- (v1);
\path[line] (v1) -- (v2) node[midway, above] {$A \land B$};
\path[line] (v2) -- (out_and);

\node[above=0.5cm of v1, xshift=1.5cm] {\textbf{Conjonction (ET) : vannes en série}};

% Nodes for OR logic (Parallel)
\node[valve, below=3cm of v1] (v3) {Vanne $A$};
\node[valve, below=5cm of v1] (v4) {Vanne $B$};
\coordinate[left=1.5cm of v3, yshift=-1cm] (split_in);
\coordinate[right=1.5cm of v3, yshift=-1cm] (split_out);
\node[left=1cm of split_in] (in_or) {Flux entrant};
\node[right=1cm of split_out] (out_or) {Flux sortant};

% Paths for OR
\path[line] (in_or) -- (split_in);
\path[draw, very thick] (split_in) |- (v3);
\path[draw, very thick] (split_in) |- (v4);
\path[draw, very thick] (v3) -| (split_out) node[pos=0.25, above] {$A \lor B$};
\path[draw, very thick] (v4) -| (split_out);
\path[line] (split_out) -- (out_or);

\node[above=0.5cm of v3, xshift=1.5cm] {\textbf{Disjonction (OU) : vannes en parallèle}};

\end{tikzpicture}
\caption{Analogie physique des connecteurs logiques à l'aide de réseaux de distribution de flux d'eau.}
\label{fig:logic-gates}
\end{figure}


Pourquoi a-t-on ressenti le besoin d'inventer un tel formalisme ? Le
langage naturel (le français, l'anglais, etc.) est parsemé d'ambiguïtés
sémantiques. Le mot ``ou'', par exemple, exprime tantôt une disjonction
exclusive (``fromage ou dessert'', mais pas les deux), tantôt une
disjonction inclusive (``recherche développeur maîtrisant le C++ ou le
Python'', de préférence les deux). De plus, l'esprit humain est sujet
aux biais cognitifs et aux sophismes.

La création d'une langue artificielle universelle, libre de toute
ambiguïté contextuelle, a permis de détacher la forme d'un raisonnement
de son contenu. C'est le projet leibnizien de la \emph{characteristica
universalis} et du \emph{calculus ratiocinator} : transformer toute
dispute philosophique ou scientifique en un simple calcul mathématique.
``Calculons !'' disait Leibniz.

Gottlob Frege, à la fin du XIXe siècle, a concrétisé ce rêve avec son
\emph{Idéographie} (\emph{Begriffsschrift}), jetant les bases de la
logique mathématique moderne.

Dans cette optique, les variables propositionnelles agissent comme des
interrupteurs élémentaires, et les formules logiques sont des montages
complexes d'interrupteurs.

Construire une table de vérité revient à dresser la carte complète des
tensions de notre circuit pour toutes les configurations possibles des
interrupteurs d'entrée.

Si un circuit produit une tension de sortie positive (Vrai) quelle que
soit la position des interrupteurs d'entrée, nous avons affaire à une
tautologie : une vérité structurelle inébranlable, indépendante du monde
réel. À l'inverse, si la sortie est toujours nulle, c'est une
contradiction.

Cette dissociation fondamentale entre la syntaxe (la forme de la
formule) et la sémantique (sa signification ou sa valeur de vérité dans
un monde donné) est la clé de voûte de la logique. Elle montre comment
des règles de manipulation purement symboliques peuvent refléter
fidèlement des réalités sémantiques.

En manipulant des symboles selon des règles strictes, sans même savoir
ce qu'ils représentent, on préserve la vérité. C'est cette nature
purement mécanique qui permettra plus tard aux ordinateurs de manipuler
des concepts abstraits et de raisonner de manière autonome.

\begin{center}\rule{0.5\linewidth}{0.5pt}\end{center}

\hypertarget{formalisation-rigueur-acaduxe9mique}{%
\subsection{2. Formalisation \& Rigueur
Académique}\label{formalisation-rigueur-acaduxe9mique}}

Pour structurer rigoureusement le calcul des propositions (ou logique
propositionnelle classique), nous devons définir précisément sa syntaxe
(l'ensemble des formules bien formées) et sa sémantique (l'évaluation de
ces formules).

\hypertarget{a.-syntaxe-du-calcul-des-propositions}{%
\subsubsection{A. Syntaxe du calcul des
propositions}\label{a.-syntaxe-du-calcul-des-propositions}}

Soit \(\mathcal{P}\) un ensemble dénombrable de symboles appelés
\textbf{variables propositionnelles} (ou atomes), généralement notés
\(p, q, r, \dots\) ou \(p_1, p_2, \dots\)

On définit l'alphabet de notre langage par la réunion de
\(\mathcal{P}\), des connecteurs logiques
\(\{\neg, \lor, \land, \Rightarrow, \Leftrightarrow\}\) et des symboles
de ponctuation \(\{\) ( , ) \(\}\).

L'ensemble \(\mathcal{F}\) des \textbf{formules bien formées} (FBF) est
défini de manière inductive comme le plus petit ensemble de mots sur cet
alphabet satisfaisant aux règles suivantes : 1. \textbf{Règle de base :}
Si \(p \in \mathcal{P}\), alors \(p \in \mathcal{F}\) (les variables
propositionnelles sont des formules atomiques). 2. \textbf{Règle
d'induction unaire :} Si \(A \in \mathcal{F}\), alors
\((\neg A) \in \mathcal{F}\). 3. \textbf{Règle d'induction binaire :} Si
\(A \in \mathcal{F}\) et \(B \in \mathcal{F}\), alors
\((A \land B) \in \mathcal{F}\), \((A \lor B) \in \mathcal{F}\),
\((A \Rightarrow B) \in \mathcal{F}\) et
\((A \Leftrightarrow B) \in \mathcal{F}\).

\emph{Remarque sur l'écriture :} Pour alléger la lecture, on omet
souvent les parenthèses extérieures et on établit des règles de priorité
des connecteurs (le connecteur \(\neg\) est plus prioritaire que
\(\land\) et \(\lor\), qui sont plus prioritaires que \(\Rightarrow\) et
\(\Leftrightarrow\)).

Le \textbf{Théorème de lecture unique} (ou d'auto-analyse) garantit que
pour toute formule bien formée non atomique \(F\), il existe un unique
connecteur principal et d'uniques sous-formules directes permettant de
la construire. Cela permet de définir des propriétés sur les formules
par induction structurelle et des fonctions par récurrence sur la
complexité des formules.

\hypertarget{b.-suxe9mantique-et-uxe9valuation}{%
\subsubsection{B. Sémantique et
Évaluation}\label{b.-suxe9mantique-et-uxe9valuation}}

La sémantique de la logique propositionnelle repose sur le principe de
bivalence (les formules sont soit vraies, soit fausses) et de
compositionnalité (la valeur de vérité d'une formule complexe dépend
uniquement de celle de ses sous-formules).

On appelle \textbf{valuation} (ou interprétation) une application
\(v : \mathcal{P} \to \{0, 1\}\). L'ensemble des valuations est noté
\(\mathcal{V} = \{0, 1\}^\mathcal{P}\).

Toute valuation \(v\) s'étend de manière unique en une application
\(\bar{v} : \mathcal{F} \to \{0, 1\}\) (que l'on notera également \(v\)
par abus de langage) définie par récurrence structurelle : - Pour tout
\(p \in \mathcal{P}\), \(\bar{v}(p) = v(p)\). -
\(\bar{v}((\neg A)) = 1 - \bar{v}(A)\). -
\(\bar{v}((A \land B)) = \min(\bar{v}(A), \bar{v}(B))\). -
\(\bar{v}((A \lor B)) = \max(\bar{v}(A), \bar{v}(B))\). -
\(\bar{v}((A \Rightarrow B)) = \max(1 - \bar{v}(A), \bar{v}(B))\). -
\(\bar{v}((A \Leftrightarrow B)) = 1 - |\bar{v}(A) - \bar{v}(B)|\).

Une formule \(A\) est dite \textbf{satisfaisable} s'il existe une
valuation \(v\) telle que \(v(A) = 1\). Dans ce cas, \(v\) est appelée
un \textbf{modèle} de \(A\) (noté \(v \models A\)). Un ensemble de
formules \(\Sigma\) est satisfaisable s'il existe une valuation qui est
modèle de toutes les formules de \(\Sigma\).

Une formule \(A\) est une \textbf{tautologie} (ou est universellement
valide) si pour toute valuation \(v\), \(v(A) = 1\). On note alors
\(\models A\). Une formule \(A\) est une \textbf{contradiction} (ou est
antilogie) si pour toute valuation \(v\), \(v(A) = 0\).

Soit \(\Sigma \subseteq \mathcal{F}\) un ensemble de formules et
\(A \in \mathcal{F}\) une formule. On dit que \(A\) est une
\textbf{conséquence sémantique} de \(\Sigma\), notée
\(\Sigma \models A\), si pour toute valuation \(v\), si \(v(B) = 1\)
pour tout \(B \in \Sigma\), alors \(v(A) = 1\).

\hypertarget{c.-systuxe8mes-formels-de-duxe9duction-syntaxe}{%
\subsubsection{C. Systèmes Formels de Déduction
(Syntaxe)}\label{c.-systuxe8mes-formels-de-duxe9duction-syntaxe}}

Parallèlement à la sémantique, on définit des systèmes de preuve
purement syntaxiques (comme les systèmes de Hilbert, la déduction
naturelle ou le calcul des séquents). Une preuve formelle d'une formule
\(A\) sous les hypothèses \(\Sigma\) est une suite finie de formules se
terminant par \(A\), où chaque formule est soit un axiome du système,
soit une formule de \(\Sigma\), soit obtenue à partir de formules
précédentes par une règle d'inférence (comme le \emph{Modus Ponens} : de
\(X\) et \(X \Rightarrow Y\), on déduit \(Y\)). On note alors
\(\Sigma \vdash A\).

Les propriétés fondamentales reliant sémantique et syntaxe sont : -
\textbf{Théorème de Correction (Soundness) :} Si \(\Sigma \vdash A\),
alors \(\Sigma \models A\). (Tout ce qui est démontrable est vrai). -
\textbf{Théorème de Complétude (Completeness) :} Si
\(\Sigma \models A\), alors \(\Sigma \vdash A\). (Tout ce qui est vrai
est démontrable).

\begin{center}\rule{0.5\linewidth}{0.5pt}\end{center}

\hypertarget{le-noyau-dur-duxe9monstrations-pas-uxe0-pas}{%
\subsection{3. Le Noyau Dur : Démonstrations
Pas-à-Pas}\label{le-noyau-dur-duxe9monstrations-pas-uxe0-pas}}

Nous présentons ici trois démonstrations fondamentales du calcul des
propositions, rédigées avec une rigueur absolue et sans aucune ellipse
mathématique.

\hypertarget{duxe9monstration-1-loi-de-de-morgan-nega-lor-b-leftrightarrow-neg-a-land-neg-b}{%
\subsubsection{\texorpdfstring{Démonstration 1 : Loi de De Morgan
\(\neg(A \lor B) \Leftrightarrow (\neg A \land \neg B)\)}{Démonstration 1 : Loi de De Morgan \textbackslash neg(A \textbackslash lor B) \textbackslash Leftrightarrow (\textbackslash neg A \textbackslash land \textbackslash neg B)}}\label{duxe9monstration-1-loi-de-de-morgan-nega-lor-b-leftrightarrow-neg-a-land-neg-b}}

Soit \(v\) une valuation quelconque. Nous voulons démontrer que pour
toutes formules \(A\) et \(B\),
\(v(\neg(A \lor B)) = v(\neg A \land \neg B)\). Nous allons procéder par
étude exhaustive des cas (valeurs de vérité de \(A\) et \(B\)).

\begin{enumerate}
\def\labelenumi{\arabic{enumi}.}
\tightlist
\item
  \textbf{Cas 1 : \(v(A) = 1\) et \(v(B) = 1\)}

  \begin{itemize}
  \tightlist
  \item
    Évaluons le membre gauche :
    \[v(A \lor B) = \max(v(A), v(B)) = \max(1, 1) = 1\]
    \[v(\neg(A \lor B)) = 1 - v(A \lor B) = 1 - 1 = 0\]
  \item
    Évaluons le membre droit : \[v(\neg A) = 1 - v(A) = 1 - 1 = 0\]
    \[v(\neg B) = 1 - v(B) = 1 - 1 = 0\]
    \[v(\neg A \land \neg B) = \min(v(\neg A), v(\neg B)) = \min(0, 0) = 0\]
  \item
    Les deux valeurs coïncident : \(0 = 0\).
  \end{itemize}
\item
  \textbf{Cas 2 : \(v(A) = 1\) et \(v(B) = 0\)}

  \begin{itemize}
  \tightlist
  \item
    Évaluons le membre gauche :
    \[v(A \lor B) = \max(v(A), v(B)) = \max(1, 0) = 1\]
    \[v(\neg(A \lor B)) = 1 - v(A \lor B) = 1 - 1 = 0\]
  \item
    Évaluons le membre droit : \[v(\neg A) = 1 - v(A) = 1 - 1 = 0\]
    \[v(\neg B) = 1 - v(B) = 1 - 0 = 1\]
    \[v(\neg A \land \neg B) = \min(v(\neg A), v(\neg B)) = \min(0, 1) = 0\]
  \item
    Les deux valeurs coïncident : \(0 = 0\).
  \end{itemize}
\item
  \textbf{Cas 3 : \(v(A) = 0\) et \(v(B) = 1\)}

  \begin{itemize}
  \tightlist
  \item
    Évaluons le membre gauche :
    \[v(A \lor B) = \max(v(A), v(B)) = \max(0, 1) = 1\]
    \[v(\neg(A \lor B)) = 1 - v(A \lor B) = 1 - 1 = 0\]
  \item
    Évaluons le membre droit : \[v(\neg A) = 1 - v(A) = 1 - 0 = 1\]
    \[v(\neg B) = 1 - v(B) = 1 - 1 = 0\]
    \[v(\neg A \land \neg B) = \min(v(\neg A), v(\neg B)) = \min(1, 0) = 0\]
  \item
    Les deux valeurs coïncident : \(0 = 0\).
  \end{itemize}
\item
  \textbf{Cas 4 : \(v(A) = 0\) et \(v(B) = 0\)}

  \begin{itemize}
  \tightlist
  \item
    Évaluons le membre gauche :
    \[v(A \lor B) = \max(v(A), v(B)) = \max(0, 0) = 0\]
    \[v(\neg(A \lor B)) = 1 - v(A \lor B) = 1 - 0 = 1\]
  \item
    Évaluons le membre droit : \[v(\neg A) = 1 - v(A) = 1 - 0 = 1\]
    \[v(\neg B) = 1 - v(B) = 1 - 0 = 1\]
    \[v(\neg A \land \neg B) = \min(v(\neg A), v(\neg B)) = \min(1, 1) = 1\]
  \item
    Les deux valeurs coïncident : \(1 = 1\).
  \end{itemize}
\end{enumerate}

Dans tous les cas possibles,
\(v(\neg(A \lor B)) = v(\neg A \land \neg B)\). La formule
\(\neg(A \lor B) \Leftrightarrow (\neg A \land \neg B)\) est donc une
tautologie.

\begin{center}\rule{0.5\linewidth}{0.5pt}\end{center}

\hypertarget{duxe9monstration-2-distributivituxe9-de-la-disjonction-par-rapport-uxe0-la-conjonction-a-lor-b-land-c-leftrightarrow-a-lor-b-land-a-lor-c}{%
\subsubsection{\texorpdfstring{Démonstration 2 : Distributivité de la
disjonction par rapport à la conjonction :
\(A \lor (B \land C) \Leftrightarrow (A \lor B) \land (A \lor C)\)}{Démonstration 2 : Distributivité de la disjonction par rapport à la conjonction : A \textbackslash lor (B \textbackslash land C) \textbackslash Leftrightarrow (A \textbackslash lor B) \textbackslash land (A \textbackslash lor C)}}\label{duxe9monstration-2-distributivituxe9-de-la-disjonction-par-rapport-uxe0-la-conjonction-a-lor-b-land-c-leftrightarrow-a-lor-b-land-a-lor-c}}

Soit \(v\) une valuation. Posons \(x = v(A)\), \(y = v(B)\) et
\(z = v(C)\), où \(x, y, z \in \{0, 1\}\). Nous devons prouver l'égalité
algébrique suivante dans l'algèbre de Boole \(\{0, 1\}\) :
\[\max(x, \min(y, z)) = \min(\max(x, y), \max(x, z))\]

Analysons les cas possibles selon la valeur de \(x\) :

\begin{enumerate}
\def\labelenumi{\arabic{enumi}.}
\tightlist
\item
  \textbf{Cas 1 : \(x = 1\)}

  \begin{itemize}
  \tightlist
  \item
    Évaluons le membre de gauche (MG) :
    \[\text{MG} = \max(1, \min(y, z))\] Puisque pour tout
    \(t \in \{0, 1\}\), \(\max(1, t) = 1\), nous avons :
    \[\text{MG} = 1\]
  \item
    Évaluons le membre de droite (MD) :
    \[\max(1, y) = 1 \quad \text{et} \quad \max(1, z) = 1\] D'où :
    \[\text{MD} = \min(1, 1) = 1\]
  \item
    Ainsi, \(\text{MG} = \text{MD}\).
  \end{itemize}
\item
  \textbf{Cas 2 : \(x = 0\)}

  \begin{itemize}
  \tightlist
  \item
    Évaluons le membre de gauche (MG) :
    \[\text{MG} = \max(0, \min(y, z))\] Puisque pour tout
    \(t \in \{0, 1\}\), \(\max(0, t) = t\), nous avons :
    \[\text{MG} = \min(y, z)\]
  \item
    Évaluons le membre de droite (MD) :
    \[\max(0, y) = y \quad \text{et} \quad \max(0, z) = z\] D'où :
    \[\text{MD} = \min(y, z)\]
  \item
    Ainsi, \(\text{MG} = \text{MD}\).
  \end{itemize}
\end{enumerate}

Dans tous les cas, l'équivalence est vérifiée, ce qui valide la
distributivité de \(\lor\) sur \(\land\).

\begin{center}\rule{0.5\linewidth}{0.5pt}\end{center}

\hypertarget{duxe9monstration-3-thuxe9oruxe8me-de-compacituxe9-de-la-logique-propositionnelle}{%
\subsubsection{Démonstration 3 : Théorème de compacité de la logique
propositionnelle}\label{duxe9monstration-3-thuxe9oruxe8me-de-compacituxe9-de-la-logique-propositionnelle}}

\begin{quote}
\textbf{Théorème :} Soit \(\Sigma \subseteq \mathcal{F}\) un ensemble
infini de formules. Si tout sous-ensemble fini
\(\Sigma_0 \subseteq \Sigma\) est satisfaisable, alors \(\Sigma\)
lui-même est satisfaisable.
\end{quote}

\textbf{Démonstration par le théorème de Tychonoff :}

Soit \(\mathcal{P}\) l'ensemble des variables propositionnelles
apparaissant dans \(\Sigma\). L'ensemble des valuations de
\(\mathcal{P}\) est l'ensemble \(\mathcal{V} = \{0, 1\}^\mathcal{P}\).
Nous pouvons munir l'ensemble discret \(\{0, 1\}\) de la topologie
discrète (qui est compacte car l'espace est fini). Par le
\textbf{Théorème de Tychonoff}, l'espace produit
\(\mathcal{V} = \{0, 1\}^\mathcal{P}\) muni de la topologie produit est
compact (espace de Cantor).

Pour chaque formule \(A \in \Sigma\), définissons l'ensemble
\(M(A) = \{v \in \mathcal{V} \mid v(A) = 1\}\), c'est-à-dire l'ensemble
des modèles de \(A\). Puisque la syntaxe d'une formule \(A\) ne fait
intervenir qu'un nombre fini de variables propositionnelles, la valeur
de vérité de \(A\) ne dépend que d'une projection finie de
\(\mathcal{V}\). Un sous-ensemble défini par un nombre fini de
coordonnées dans une topologie produit de topologies discrètes est fermé
(et même ouvert). Donc, \(M(A)\) est un fermé de l'espace compact
\(\mathcal{V}\).

L'hypothèse selon laquelle tout sous-ensemble fini
\(\Sigma_0 \subseteq \Sigma\) est satisfaisable se traduit
sémantiquement par : \[\bigcap_{A \in \Sigma_0} M(A) \neq \emptyset\]

Considérons la famille de fermés de l'espace compact \(\mathcal{V}\)
définie par \(\mathcal{C} = \{ M(A) \mid A \in \Sigma \}\). D'après le
résultat ci-dessus, cette famille possède la \textbf{propriété de
l'intersection finie} (toute intersection d'un nombre fini de membres de
la famille est non vide).

Par caractérisation de la compacité, une famille de fermés d'un espace
compact ayant la propriété de l'intersection finie a une intersection
totale non vide : \[\bigcap_{A \in \Sigma} M(A) \neq \emptyset\]

Il existe donc une valuation \(v^* \in \bigcap_{A \in \Sigma} M(A)\).
Cette valuation \(v^*\) est modèle de toutes les formules de \(\Sigma\).
Par conséquent, \(\Sigma\) est satisfaisable. \(\blacksquare\)

\begin{center}\rule{0.5\linewidth}{0.5pt}\end{center}

\hypertarget{exercices-dapplication-pratique-de-concours}{%
\subsection{4. Exercices d'Application \& Pratique de
Concours}\label{exercices-dapplication-pratique-de-concours}}

\hypertarget{exercice-1-preuve-de-la-loi-de-peirce}{%
\subsubsection{Exercice 1 : Preuve de la Loi de
Peirce}\label{exercice-1-preuve-de-la-loi-de-peirce}}

\textbf{Énoncé :} Démontrer algébriquement, puis à l'aide de
l'évaluation sémantique, que la formule
\(((A \Rightarrow B) \Rightarrow A) \Rightarrow A\) (Loi de Peirce) est
une tautologie.

\textbf{Correction pas-à-pas :}

\begin{enumerate}
\def\labelenumi{\arabic{enumi}.}
\item
  \textbf{Méthode sémantique :} Soit \(v\) une valuation. Posons
  \(a = v(A)\) et \(b = v(B)\). Nous voulons évaluer
  \(v(((A \Rightarrow B) \Rightarrow A) \Rightarrow A)\). Rappelons que
  \(v(X \Rightarrow Y) = 1\) ssi \(v(X) \le v(Y)\), ce qui s'écrit
  algébriquement \(v(X \Rightarrow Y) = \max(1 - v(X), v(Y))\).

  \begin{itemize}
  \item
    \textbf{Cas 1 : \(a = 1\).} Le membre final de notre implication
    principale est la formule \(A\). Or sa valeur est \(a = 1\). Une
    implication de la forme \(X \Rightarrow Y\) prend la valeur \(1\) si
    la conséquence \(Y\) vaut \(1\), car \(\max(1-v(X), 1) = 1\).
    Puisque la conséquence de notre implication principale est \(A\) et
    que \(v(A) = 1\), la formule entière prend la valeur \(1\).
  \item
    \textbf{Cas 2 : \(a = 0\).}

    \begin{itemize}
    \tightlist
    \item
      Évaluons d'abord l'implication la plus interne :
      \[v(A \Rightarrow B) = \max(1 - a, b) = \max(1 - 0, b) = \max(1, b) = 1\]
    \item
      Évaluons le membre gauche de l'implication principale :
      \[v((A \Rightarrow B) \Rightarrow A) = \max(1 - v(A \Rightarrow B), a) = \max(1 - 1, 0) = \max(0, 0) = 0\]
    \item
      Évaluons enfin l'implication principale :
      \[v(((A \Rightarrow B) \Rightarrow A) \Rightarrow A) = \max(1 - v((A \Rightarrow B) \Rightarrow A), a) = \max(1 - 0, 0) = \max(1, 0) = 1\]
    \end{itemize}
  \end{itemize}

  Dans les deux cas possibles, la valeur de vérité de la formule est
  \(1\). C'est donc une tautologie.
\item
  \textbf{Méthode algébrique (Simplification propositionnelle) :}
  Utilisons l'équivalence \(X \Rightarrow Y \equiv \neg X \lor Y\).

  \begin{itemize}
  \tightlist
  \item
    Étape 1 : Réécriture de l'implication interne :
    \[(A \Rightarrow B) \Rightarrow A \equiv \neg(\neg A \lor B) \lor A\]
  \item
    Étape 2 : Application de la loi de De Morgan sur la négation externe
    :
    \[\neg(\neg A \lor B) \lor A \equiv (\neg\neg A \land \neg B) \lor A \equiv (A \land \neg B) \lor A\]
  \item
    Étape 3 : Par absorption dans l'algèbre de Boole,
    \((A \land \neg B) \lor A \equiv A\). \emph{Justification pas-à-pas
    de l'absorption :}
    \[(A \land \neg B) \lor A \equiv (A \land \neg B) \lor (A \land 1) \equiv A \land (\neg B \lor 1) \equiv A \land 1 \equiv A\]
  \item
    Étape 4 : Substituons ce résultat dans la formule complète :
    \[((A \Rightarrow B) \Rightarrow A) \Rightarrow A \equiv A \Rightarrow A\]
  \item
    Étape 5 : Simplifions \(A \Rightarrow A\) :
    \[A \Rightarrow A \equiv \neg A \lor A \equiv 1\] La formule est
    bien équivalente à la constante Vrai, c'est une tautologie.
  \end{itemize}
\end{enumerate}

\begin{center}\rule{0.5\linewidth}{0.5pt}\end{center}

\hypertarget{exercice-2-compluxe9tude-fonctionnelle-du-connecteur-de-sheffer-nand}{%
\subsubsection{Exercice 2 : Complétude fonctionnelle du connecteur de
Sheffer
(NAND)}\label{exercice-2-compluxe9tude-fonctionnelle-du-connecteur-de-sheffer-nand}}

\textbf{Énoncé :} On définit le connecteur binaire \(\mid\) (appelé
barre de Sheffer ou NAND) par la table de vérité suivante :
\(v(A \mid B) = 1 - \min(v(A), v(B))\). Démontrer que le singleton
\(\{\mid\}\) est un système complet de connecteurs (c'est-à-dire que
n'importe quelle formule propositionnelle peut s'écrire uniquement à
l'aide de ce connecteur).

\textbf{Correction pas-à-pas :}

\begin{enumerate}
\def\labelenumi{\arabic{enumi}.}
\item
  \textbf{Rappel théorique :} On sait que le couple de connecteurs
  \(\{\neg, \land\}\) est fonctionnellement complet. Pour montrer que
  \(\{\mid\}\) est complet, il suffit d'exprimer les opérations de
  négation (\(\neg\)) et de conjonction (\(\land\)) en utilisant
  exclusivement le connecteur \(\mid\).
\item
  \textbf{Expression de la négation \(\neg A\) :} Montrons que
  \(\neg A\) est équivalent à \(A \mid A\). Soit \(v\) une valuation
  quelconque. \[v(A \mid A) = 1 - \min(v(A), v(A)) = 1 - v(A)\] Or,
  \(v(\neg A) = 1 - v(A)\). Les deux évaluations coïncident. Nous avons
  donc : \[\neg A \equiv A \mid A\]
\item
  \textbf{Expression de la conjonction \(A \land B\) :} Montrons que
  \(A \land B\) est équivalent à \((A \mid B) \mid (A \mid B)\). Soit
  \(v\) une valuation. Posons \(X = A \mid B\). Par la définition de la
  négation établie au point précédent : \[v(X \mid X) = 1 - v(X)\]
  Remplaçons \(v(X)\) par sa définition :
  \[v(X) = 1 - \min(v(A), v(B)) = v(\neg(A \land B))\] D'où :
  \[v(X \mid X) = 1 - (1 - \min(v(A), v(B))) = \min(v(A), v(B)) = v(A \land B)\]
  Ainsi, les formules coïncident. Nous avons établi :
  \[A \land B \equiv (A \mid B) \mid (A \mid B)\]
\item
  \textbf{Conclusion :} Puisque nous pouvons traduire la négation
  (\(\neg\)) et la conjonction (\(\land\)) uniquement à l'aide de la
  barre de Sheffer, et que le système \(\{\neg, \land\}\) est complet,
  le singleton \(\{\mid\}\) est à son tour un système complet de
  connecteurs.
\end{enumerate}

\begin{center}\rule{0.5\linewidth}{0.5pt}\end{center}

\hypertarget{ancrage-application-en-intelligence-artificielle}{%
\subsection{5. Ancrage \& Application en Intelligence
Artificielle}\label{ancrage-application-en-intelligence-artificielle}}

La logique propositionnelle, bien que formulée au XIXe siècle, est au
cœur des révolutions technologiques de l'Intelligence Artificielle
contemporaine, qu'elle soit symbolique ou hybride.

\hypertarget{le-sat-solving-le-moteur-de-lia-symbolique}{%
\subsubsection{Le SAT-Solving : Le moteur de l'IA
symbolique}\label{le-sat-solving-le-moteur-de-lia-symbolique}}

Le problème de la décidabilité de la satisfaisabilité d'une formule
propositionnelle (problème \textbf{SAT}) est le premier problème à avoir
été prouvé \textbf{NP-complet} par Stephen Cook (1971) et Leonid Levin.
Malgré cette complexité théorique redoutable en pire des cas, les
chercheurs en IA ont développé des moteurs de résolution de contraintes
logiques d'une efficacité spectaculaire (les \textbf{SAT-Solvers}),
capables de traiter des formules comportant des millions de variables et
de clauses.

Le cœur de ces moteurs repose sur l'algorithme historique \textbf{DPLL}
(Davis-Putnam-Logemann-Loveland) et sa version moderne \textbf{CDCL}
(Conflict-Driven Clause Learning). L'algorithme opère sur des formules
écrites sous Forme Normale Conjonctive (CNF) -- c'est-à-dire des
conjonctions de clauses, où chaque clause est une disjonction de
littéraux (variables ou négation de variables). Le principe du CDCL est
le suivant : 1. \textbf{Assignation de variables (décisions) :} On
choisit une variable et on lui affecte arbitrairement une valeur (\(0\)
ou \(1\)). 2. \textbf{Propagation unitaire (Boolean Constraint
Propagation - BCP) :} Si une clause ne contient plus qu'un littéral non
assigné, cette variable doit être forcée à la valeur qui rend le
littéral vrai. 3. \textbf{Analyse de conflit et Apprentissage :} Si une
contradiction (conflit) émerge (une clause devenant entièrement fausse),
le solveur n'effectue pas un simple retour en arrière classique
(\emph{backtracking}). Il analyse le graphe des implications qui a mené
au conflit, extrait une clause dite de conflit (la cause minimale de
l'erreur) et l'ajoute à sa base de connaissances pour ne plus jamais
reproduire la même suite d'erreurs. 4. \textbf{Retour arrière
non-chronologique (\emph{backjumping}) :} Le solveur remonte directement
au niveau de décision qui a provoqué le conflit.

Ces moteurs sont utilisés en IA pour : - \textbf{La vérification
formelle de modèles :} S'assurer que le code d'un pilote automatique ou
d'un système critique (médical, aérospatial) ne peut jamais se retrouver
dans un état instable ou interdit (modélisé par une contradiction
logique). - \textbf{La planification automatisée :} Trouver une séquence
d'actions permettant à un robot d'atteindre un but sous des contraintes
logiques strictes (l'approche SATPlan). - \textbf{La vérification
formelle des Réseaux de Neurones :} Des outils récents encodent le
comportement local de neurones avec des fonctions d'activation linéaires
par morceaux (ReLU) sous forme de contraintes logiques pour prouver
qu'un réseau profond ne subira pas de failles de type ``exemples
contradictoires'' ou ``exemples antagonistes'' (\emph{adversarial
examples}).

\begin{center}\rule{0.5\linewidth}{0.5pt}\end{center}


\begin{figure}[htbp]
\centering
\begin{tikzpicture}[
    level 1/.style={sibling distance=4cm},
    level 2/.style={sibling distance=2cm},
    node/.style={circle, draw=red!60, fill=red!5, very thick, minimum size=8mm},
    leaf/.style={rectangle, draw=green!60, fill=green!5, very thick, minimum size=8mm},
    conflict/.style={rectangle, draw=orange!80, fill=orange!10, very thick, minimum size=8mm},
    edge/.style={draw, very thick, -latex'}
]

\node[node] (p) {$P$}
    child { node[node] (q1) {$Q$}
        child { node[leaf] (sat) {SAT} edge from parent node[left, xshift=-0.1cm] {$0$} }
        child { node[conflict] (c1) {Conflit} edge from parent node[right, xshift=0.1cm] {$1$} }
        edge from parent node[left, xshift=-0.1cm] {$0$}
    }
    child { node[node] (q2) {$Q$}
        child { node[conflict] (c2) {Conflit} edge from parent node[left, xshift=-0.1cm] {$0$} }
        child { node[conflict] (c3) {Conflit} edge from parent node[right, xshift=0.1cm] {$1$} }
        edge from parent node[right, xshift=0.1cm] {$1$}
    };

\end{tikzpicture}
\caption{Arbre de recherche (Backtracking/DPLL) explorant l'espace des valuations d'une formule propositionnelle.}
\label{fig:dpll-tree}
\end{figure}


\hypertarget{liens-suxe9mantiques-maillage-obsidian}{%
\subsection{6. Liens Sémantiques \& Maillage
Obsidian}\label{liens-suxe9mantiques-maillage-obsidian}}

\begin{itemize}
\tightlist
\item
  \textbf{Concepts Précédents requis :} Aucun (Jalon initial)
\item
  \textbf{Concepts Futurs dépendants :} {[}{[}Jalon 2 (Méthodes de
  raisonnement){]}{]}, {[}{[}Jalon 3 (Quantification){]}{]}, {[}{[}Jalon
  12 (Livrable IA){]}{]}, {[}{[}Jalon 133 (Modele PAC){]}{]}
\end{itemize}


\part{Exercices d'application et pratique de concours}
\hypertarget{exercice-1-construction-et-interpruxe9tation-de-tables-de-vuxe9rituxe9}{%
\section{Exercice 1 : Construction et interprétation de tables de
vérité}\label{exercice-1-construction-et-interpruxe9tation-de-tables-de-vuxe9rituxe9}}

\hypertarget{uxe9noncuxe9}{%
\subsection{Énoncé}\label{uxe9noncuxe9}}

Soient \(P\), \(Q\) et \(R\) trois variables propositionnelles
indépendantes. On considère la formule propositionnelle complexe
suivante, notée \(F\) :
\[F = \big( (P \Rightarrow Q) \land (Q \Rightarrow R) \big) \Rightarrow (P \Rightarrow R)\]

\begin{enumerate}
\def\labelenumi{\arabic{enumi}.}
\tightlist
\item
  Expliquer pourquoi le nombre de lignes de la table de vérité pour
  cette formule est égal à \(8\).
\item
  Construire pas-à-pas la table de vérité complète de la formule \(F\),
  en détaillant chaque colonne intermédiaire.
\item
  Déduire de la table de vérité la nature de la formule \(F\)
  (satisfaisable, tautologie ou contradiction).
\end{enumerate}

\begin{center}\rule{0.5\linewidth}{0.5pt}\end{center}

\hypertarget{correction-duxe9tailluxe9e}{%
\subsection{Correction Détaillée}\label{correction-duxe9tailluxe9e}}

\hypertarget{question-1-nombre-de-lignes-de-la-table-de-vuxe9rituxe9}{%
\subsubsection{Question 1 : Nombre de lignes de la table de
vérité}\label{question-1-nombre-de-lignes-de-la-table-de-vuxe9rituxe9}}

Chaque variable propositionnelle (\(P\), \(Q\), \(R\)) peut prendre de
manière indépendante deux valeurs de vérité distinctes : \(1\) (Vrai) ou
\(0\) (Faux). Puisqu'il y a \(3\) variables propositionnelles
distinctes, l'ensemble de toutes les configurations possibles
(valuations) correspond au produit cartésien \(\{0, 1\}^3\). Le nombre
cardinal de cet ensemble est : \[\text{Cardinal}(\{0, 1\}^3) = 2^3 = 8\]
Il y a donc exactement \(8\) valuations distinctes à analyser, ce qui se
traduit par \(8\) lignes dans notre table de vérité.

\hypertarget{question-2-construction-de-la-table-de-vuxe9rituxe9}{%
\subsubsection{Question 2 : Construction de la table de
vérité}\label{question-2-construction-de-la-table-de-vuxe9rituxe9}}

Nous allons définir les colonnes de notre tableau dans l'ordre
d'évaluation des sous-formules : - Colonnes 1, 2, 3 : Variables de base
\(P\), \(Q\), \(R\). - Colonne 4 : \(P \Rightarrow Q\) (vaut \(0\)
uniquement si \(P=1\) et \(Q=0\)). - Colonne 5 : \(Q \Rightarrow R\)
(vaut \(0\) uniquement si \(Q=1\) et \(R=0\)). - Colonne 6 :
\((P \Rightarrow Q) \land (Q \Rightarrow R)\) (conjonction des colonnes
4 et 5 ; vaut \(1\) si et seulement si les deux valent \(1\)). - Colonne
7 : \(P \Rightarrow R\) (vaut \(0\) uniquement si \(P=1\) et \(R=0\)). -
Colonne 8 : \(F = \text{Col } 6 \Rightarrow \text{Col } 7\) (vaut \(0\)
uniquement si la colonne 6 vaut \(1\) et la colonne 7 vaut \(0\)).

Dressons le tableau exhaustif :

\begin{longtable}[]{@{}
  >{\centering\arraybackslash}p{(\columnwidth - 16\tabcolsep) * \real{0.0450}}
  >{\centering\arraybackslash}p{(\columnwidth - 16\tabcolsep) * \real{0.0270}}
  >{\centering\arraybackslash}p{(\columnwidth - 16\tabcolsep) * \real{0.0270}}
  >{\centering\arraybackslash}p{(\columnwidth - 16\tabcolsep) * \real{0.0270}}
  >{\centering\arraybackslash}p{(\columnwidth - 16\tabcolsep) * \real{0.1532}}
  >{\centering\arraybackslash}p{(\columnwidth - 16\tabcolsep) * \real{0.1532}}
  >{\centering\arraybackslash}p{(\columnwidth - 16\tabcolsep) * \real{0.3874}}
  >{\centering\arraybackslash}p{(\columnwidth - 16\tabcolsep) * \real{0.1532}}
  >{\centering\arraybackslash}p{(\columnwidth - 16\tabcolsep) * \real{0.0270}}@{}}
\toprule\noalign{}
\begin{minipage}[b]{\linewidth}\centering
Ligne
\end{minipage} & \begin{minipage}[b]{\linewidth}\centering
\(P\)
\end{minipage} & \begin{minipage}[b]{\linewidth}\centering
\(Q\)
\end{minipage} & \begin{minipage}[b]{\linewidth}\centering
\(R\)
\end{minipage} & \begin{minipage}[b]{\linewidth}\centering
\(P \Rightarrow Q\)
\end{minipage} & \begin{minipage}[b]{\linewidth}\centering
\(Q \Rightarrow R\)
\end{minipage} & \begin{minipage}[b]{\linewidth}\centering
\((P \Rightarrow Q) \land (Q \Rightarrow R)\)
\end{minipage} & \begin{minipage}[b]{\linewidth}\centering
\(P \Rightarrow R\)
\end{minipage} & \begin{minipage}[b]{\linewidth}\centering
\(F\)
\end{minipage} \\
\midrule\noalign{}
\endhead
\bottomrule\noalign{}
\endlastfoot
1 & 1 & 1 & 1 & 1 & 1 & 1 & 1 & \textbf{1} \\
2 & 1 & 1 & 0 & 1 & 0 & 0 & 0 & \textbf{1} \\
3 & 1 & 0 & 1 & 0 & 1 & 0 & 1 & \textbf{1} \\
4 & 1 & 0 & 0 & 0 & 1 & 0 & 0 & \textbf{1} \\
5 & 0 & 1 & 1 & 1 & 1 & 1 & 1 & \textbf{1} \\
6 & 0 & 1 & 0 & 1 & 0 & 0 & 1 & \textbf{1} \\
7 & 0 & 0 & 1 & 1 & 1 & 1 & 1 & \textbf{1} \\
8 & 0 & 0 & 0 & 1 & 1 & 1 & 1 & \textbf{1} \\
\end{longtable}

Vérifions pas-à-pas la ligne 2 : - \(P=1, Q=1, R=0\). -
\(P \Rightarrow Q = 1 \Rightarrow 1 = 1\). -
\(Q \Rightarrow R = 1 \Rightarrow 0 = 0\). - Conjonction :
\(1 \land 0 = 0\). - \(P \Rightarrow R = 1 \Rightarrow 0 = 0\). -
Implication principale : \(0 \Rightarrow 0 = 1\). La ligne 2 est
correcte.

Vérifions pas-à-pas la ligne 6 : - \(P=0, Q=1, R=0\). -
\(P \Rightarrow Q = 0 \Rightarrow 1 = 1\). -
\(Q \Rightarrow R = 1 \Rightarrow 0 = 0\). - Conjonction :
\(1 \land 0 = 0\). - \(P \Rightarrow R = 0 \Rightarrow 0 = 1\). -
Implication principale : \(0 \Rightarrow 1 = 1\). La ligne 6 est
correcte.

\hypertarget{question-3-analyse-du-ruxe9sultat}{%
\subsubsection{Question 3 : Analyse du
résultat}\label{question-3-analyse-du-ruxe9sultat}}

La colonne finale associée à la formule \(F\) ne contient que des
valeurs de vérité \(1\) (Vrai) pour l'ensemble des \(8\) lignes. Cela
signifie que pour toute valuation \(v \in \{0, 1\}^3\), nous avons :
\[v(F) = 1\] Par définition, la formule \(F\) est donc une
\textbf{tautologie} (notée \(\models F\)). Elle est sémantiquement
universellement valide. De plus, étant toujours vraie, elle est
également satisfaisable. Il s'agit de la loi de transitivité de
l'implication logique (ou syllogisme hypothétique).
\n
\hypertarget{exercice-2-simplification-suxe9mantique-par-lalguxe8bre-de-boole}{%
\section{Exercice 2 : Simplification sémantique par l'algèbre de
Boole}\label{exercice-2-simplification-suxe9mantique-par-lalguxe8bre-de-boole}}

\hypertarget{uxe9noncuxe9}{%
\subsection{Énoncé}\label{uxe9noncuxe9}}

Soient \(P\) et \(Q\) deux variables propositionnelles. On considère la
formule logique suivante :
\[A = \big( P \lor (\neg P \land Q) \big) \land \big( \neg Q \lor (P \land \neg Q) \big)\]

\begin{enumerate}
\def\labelenumi{\arabic{enumi}.}
\tightlist
\item
  Simplifier algébriquement le membre de gauche
  \(L = P \lor (\neg P \land Q)\) en utilisant les axiomes de
  distributivité et d'absorption.
\item
  Simplifier algébriquement le membre de droite
  \(R = \neg Q \lor (P \land \neg Q)\) en utilisant la loi d'absorption.
\item
  En déduire la forme simplifiée finale de la formule \(A\).
\item
  Démontrer le résultat obtenu en dressant la table de vérité de la
  formule originale \(A\).
\end{enumerate}

\begin{center}\rule{0.5\linewidth}{0.5pt}\end{center}

\hypertarget{correction-duxe9tailluxe9e}{%
\subsection{Correction Détaillée}\label{correction-duxe9tailluxe9e}}

\hypertarget{question-1-simplification-du-membre-de-gauche-l-p-lor-neg-p-land-q}{%
\subsubsection{\texorpdfstring{Question 1 : Simplification du membre de
gauche
\(L = P \lor (\neg P \land Q)\)}{Question 1 : Simplification du membre de gauche L = P \textbackslash lor (\textbackslash neg P \textbackslash land Q)}}\label{question-1-simplification-du-membre-de-gauche-l-p-lor-neg-p-land-q}}

Appliquons la règle de \textbf{distributivité de la disjonction
(\(\lor\)) par rapport à la conjonction (\(\land\))} :
\[L = (P \lor \neg P) \land (P \lor Q)\]

Or, d'après la loi du tiers exclu (complémentarité dans les algèbres de
Boole) : \[P \lor \neg P \equiv 1\]

Substituons cette valeur dans l'expression :
\[L \equiv 1 \land (P \lor Q)\]

Puisque \(1\) (Vrai) est l'élément neutre pour la conjonction
(\(\land\)) : \[L \equiv P \lor Q\]

\begin{center}\rule{0.5\linewidth}{0.5pt}\end{center}

\hypertarget{question-2-simplification-du-membre-de-droite-r-neg-q-lor-p-land-neg-q}{%
\subsubsection{\texorpdfstring{Question 2 : Simplification du membre de
droite
\(R = \neg Q \lor (P \land \neg Q)\)}{Question 2 : Simplification du membre de droite R = \textbackslash neg Q \textbackslash lor (P \textbackslash land \textbackslash neg Q)}}\label{question-2-simplification-du-membre-de-droite-r-neg-q-lor-p-land-neg-q}}

Nous pouvons appliquer directement la \textbf{loi d'absorption} de
l'algèbre de Boole. Rappelons la loi d'absorption générale : pour tous
éléments \(x\) et \(y\), \(x \lor (y \land x) \equiv x\).

Ici, posons \(x = \neg Q\) et \(y = P\). L'expression devient :
\[R = x \lor (y \land x)\]

Par application directe de la loi d'absorption :
\[R \equiv x \equiv \neg Q\]

\begin{center}\rule{0.5\linewidth}{0.5pt}\end{center}

\hypertarget{question-3-simplification-finale-de-la-formule-a}{%
\subsubsection{\texorpdfstring{Question 3 : Simplification finale de la
formule
\(A\)}{Question 3 : Simplification finale de la formule A}}\label{question-3-simplification-finale-de-la-formule-a}}

La formule complète s'écrit \(A = L \land R\). Substituons les
expressions simplifiées obtenues aux questions 1 et 2 :
\[A \equiv (P \lor Q) \land \neg Q\]

Appliquons la distributivité de la conjonction (\(\land\)) par rapport à
la disjonction (\(\lor\)) :
\[A \equiv (P \land \neg Q) \lor (Q \land \neg Q)\]

D'après le principe de non-contradiction, \(Q \land \neg Q \equiv 0\).
L'expression devient : \[A \equiv (P \land \neg Q) \lor 0\]

Puisque \(0\) (Faux) est l'élément neutre pour la disjonction (\(\lor\))
: \[A \equiv P \land \neg Q\]

La formule originale complexe \(A\) est donc sémantiquement équivalente
à la conjonction simple de \(P\) et de la négation de \(Q\).

\begin{center}\rule{0.5\linewidth}{0.5pt}\end{center}

\hypertarget{question-4-validation-par-table-de-vuxe9rituxe9}{%
\subsubsection{Question 4 : Validation par table de
vérité}\label{question-4-validation-par-table-de-vuxe9rituxe9}}

Construisons la table de vérité pour vérifier l'équivalence
\(A \equiv P \land \neg Q\).

\begin{longtable}[]{@{}
  >{\centering\arraybackslash}p{(\columnwidth - 18\tabcolsep) * \real{0.0192}}
  >{\centering\arraybackslash}p{(\columnwidth - 18\tabcolsep) * \real{0.0192}}
  >{\centering\arraybackslash}p{(\columnwidth - 18\tabcolsep) * \real{0.0513}}
  >{\centering\arraybackslash}p{(\columnwidth - 18\tabcolsep) * \real{0.0513}}
  >{\centering\arraybackslash}p{(\columnwidth - 18\tabcolsep) * \real{0.1026}}
  >{\centering\arraybackslash}p{(\columnwidth - 18\tabcolsep) * \real{0.1859}}
  >{\centering\arraybackslash}p{(\columnwidth - 18\tabcolsep) * \real{0.1026}}
  >{\centering\arraybackslash}p{(\columnwidth - 18\tabcolsep) * \real{0.2179}}
  >{\centering\arraybackslash}p{(\columnwidth - 18\tabcolsep) * \real{0.0962}}
  >{\centering\arraybackslash}p{(\columnwidth - 18\tabcolsep) * \real{0.1538}}@{}}
\toprule\noalign{}
\begin{minipage}[b]{\linewidth}\centering
\(P\)
\end{minipage} & \begin{minipage}[b]{\linewidth}\centering
\(Q\)
\end{minipage} & \begin{minipage}[b]{\linewidth}\centering
\(\neg P\)
\end{minipage} & \begin{minipage}[b]{\linewidth}\centering
\(\neg Q\)
\end{minipage} & \begin{minipage}[b]{\linewidth}\centering
\(\neg P \land Q\)
\end{minipage} & \begin{minipage}[b]{\linewidth}\centering
\(L = P \lor (\neg P \land Q)\)
\end{minipage} & \begin{minipage}[b]{\linewidth}\centering
\(P \land \neg Q\)
\end{minipage} & \begin{minipage}[b]{\linewidth}\centering
\(R = \neg Q \lor (P \land \neg Q)\)
\end{minipage} & \begin{minipage}[b]{\linewidth}\centering
\(A = L \land R\)
\end{minipage} & \begin{minipage}[b]{\linewidth}\centering
\(P \land \neg Q\) (cible)
\end{minipage} \\
\midrule\noalign{}
\endhead
\bottomrule\noalign{}
\endlastfoot
1 & 1 & 0 & 0 & 0 & 1 & 0 & 0 & \textbf{0} & \textbf{0} \\
1 & 0 & 0 & 1 & 0 & 1 & 1 & 1 & \textbf{1} & \textbf{1} \\
0 & 1 & 1 & 0 & 1 & 1 & 0 & 0 & \textbf{0} & \textbf{0} \\
0 & 0 & 1 & 1 & 0 & 0 & 0 & 1 & \textbf{0} & \textbf{0} \\
\end{longtable}

Les colonnes \(A = L \land R\) (colonne 9) et \(P \land \neg Q\)
(colonne 10) sont parfaitement identiques pour toutes les lignes
\(\{0, 1, 0, 0\}\). L'équivalence logique \(A \equiv P \land \neg Q\)
est donc rigoureusement démontrée par la méthode sémantique.
\n
\hypertarget{exercice-3-dualituxe9-des-connecteurs-et-foncteurs-de-vuxe9rituxe9}{%
\section{Exercice 3 : Dualité des connecteurs et foncteurs de
vérité}\label{exercice-3-dualituxe9-des-connecteurs-et-foncteurs-de-vuxe9rituxe9}}

\hypertarget{uxe9noncuxe9}{%
\subsection{Énoncé}\label{uxe9noncuxe9}}

On restreint dans cet exercice notre langage de formules à l'ensemble
des formules positives construites exclusivement à partir des variables
de \(\mathcal{P}\), des connecteurs de conjonction \(\land\), de
disjonction \(\lor\) et de négation \(\neg\).

Soit \(F\) une formule. On définit sa \textbf{formule duale}, notée
\(F^*\), par récurrence structurelle sur la complexité de \(F\) comme
suit : 1. Si \(F = p\) (avec \(p \in \mathcal{P}\)), alors \(F^* = p\).
2. Si \(F = \neg A\), alors \(F^* = \neg (A^*)\). 3. Si
\(F = A \land B\), alors \(F^* = A^* \lor B^*\). 4. Si \(F = A \lor B\),
alors \(F^* = A^* \land B^*\).

Pour toute valuation \(v \in \{0, 1\}^\mathcal{P}\), on définit la
valuation conjuguée \(\bar{v}\) par \(\bar{v}(p) = 1 - v(p)\) pour tout
\(p \in \mathcal{P}\).

\begin{enumerate}
\def\labelenumi{\arabic{enumi}.}
\tightlist
\item
  Démontrer par induction structurelle sur la formule \(F\) le théorème
  de dualité sémantique suivant :
  \[\text{Pour toute valuation } v, \quad v(F^*) = 1 - \bar{v}(F)\]
\item
  En déduire que \(F\) est une tautologie si et seulement si la formule
  \(\neg (F^*[\neg p_1/p_1, \dots, \neg p_n/p_n])\) en est une, où la
  notation \(G[\neg p/p]\) représente la substitution uniforme de la
  variable \(p\) par son opposée \(\neg p\).
\item
  Déterminer et simplifier la duale de la formule :
  \[G = (P \land Q) \lor R\]
\end{enumerate}

\begin{center}\rule{0.5\linewidth}{0.5pt}\end{center}

\hypertarget{correction-duxe9tailluxe9e}{%
\subsection{Correction Détaillée}\label{correction-duxe9tailluxe9e}}

\hypertarget{question-1-preuve-du-thuxe9oruxe8me-de-dualituxe9-suxe9mantique}{%
\subsubsection{Question 1 : Preuve du théorème de dualité
sémantique}\label{question-1-preuve-du-thuxe9oruxe8me-de-dualituxe9-suxe9mantique}}

Nous allons procéder par induction structurelle sur la formule \(F\).
Soit la propriété \(\mathcal{H}(F)\) : « Pour toute valuation \(v\),
\(v(F^*) = 1 - \bar{v}(F)\) ».

\hypertarget{uxe9tape-de-base-formules-atomiques}{%
\paragraph{Étape de base (formules atomiques)
:}\label{uxe9tape-de-base-formules-atomiques}}

Soit \(F = p\) avec \(p \in \mathcal{P}\). - D'après la règle 1 de
définition du dual, \(F^* = p\). Donc \(v(F^*) = v(p)\). - Évaluons le
membre droit : \[1 - \bar{v}(F) = 1 - \bar{v}(p)\] Par définition de la
valuation conjuguée, \(\bar{v}(p) = 1 - v(p)\). D'où :
\[1 - \bar{v}(F) = 1 - (1 - v(p)) = v(p)\] Les deux membres coïncident.
L'étape de base est vérifiée.

\hypertarget{uxe9tape-dinduction}{%
\paragraph{Étape d'induction :}\label{uxe9tape-dinduction}}

Supposons que \(\mathcal{H}(A)\) et \(\mathcal{H}(B)\) sont vraies pour
deux formules \(A\) et \(B\).

\begin{itemize}
\tightlist
\item
  \textbf{Sous-cas 1 : \(F = \neg A\)} D'après la définition du dual,
  \(F^* = \neg(A^*)\).

  \begin{itemize}
  \tightlist
  \item
    Évaluons \(v(F^*)\) : \[v(F^*) = v(\neg (A^*)) = 1 - v(A^*)\] Par
    hypothèse de récurrence sémantique sur \(A\), nous avons
    \(v(A^*) = 1 - \bar{v}(A)\). D'où :
    \[v(F^*) = 1 - (1 - \bar{v}(A)) = \bar{v}(A)\]
  \item
    Évaluons le membre droit \(1 - \bar{v}(F)\) :
    \[1 - \bar{v}(\neg A) = 1 - (1 - \bar{v}(A)) = \bar{v}(A)\] Les deux
    valeurs coïncident. \(\mathcal{H}(\neg A)\) est vraie.
  \end{itemize}
\item
  \textbf{Sous-cas 2 : \(F = A \land B\)} D'après la définition,
  \(F^* = A^* \lor B^*\).

  \begin{itemize}
  \tightlist
  \item
    Évaluons \(v(F^*)\) :
    \[v(F^*) = v(A^* \lor B^*) = \max(v(A^*), v(B^*))\] En utilisant
    l'hypothèse de récurrence sémantique sur \(A\) et \(B\) :
    \[v(F^*) = \max(1 - \bar{v}(A), 1 - \bar{v}(B))\] Puisque pour tous
    réels \(a\) et \(b\), \(\max(1-a, 1-b) = 1 - \min(a, b)\), on a :
    \[v(F^*) = 1 - \min(\bar{v}(A), \bar{v}(B))\] Par définition
    sémantique de la conjonction,
    \(\min(\bar{v}(A), \bar{v}(B)) = \bar{v}(A \land B) = \bar{v}(F)\).
    D'où : \[v(F^*) = 1 - \bar{v}(F)\] Les deux valeurs coïncident.
    \(\mathcal{H}(A \land B)\) est vraie.
  \end{itemize}
\item
  \textbf{Sous-cas 3 : \(F = A \lor B\)} D'après la définition,
  \(F^* = A^* \land B^*\).

  \begin{itemize}
  \tightlist
  \item
    Évaluons \(v(F^*)\) :
    \[v(F^*) = v(A^* \land B^*) = \min(v(A^*), v(B^*))\] En utilisant
    l'hypothèse de récurrence :
    \[v(F^*) = \min(1 - \bar{v}(A), 1 - \bar{v}(B))\] Puisque
    \(\min(1-a, 1-b) = 1 - \max(a, b)\), on a :
    \[v(F^*) = 1 - \max(\bar{v}(A), \bar{v}(B)) = 1 - \bar{v}(A \lor B) = 1 - \bar{v}(F)\]
    Les deux valeurs coïncident. \(\mathcal{H}(A \lor B)\) est vraie.
  \end{itemize}
\end{itemize}

\hypertarget{conclusion-de-linduction}{%
\paragraph{Conclusion de l'induction :}\label{conclusion-de-linduction}}

La propriété \(\mathcal{H}(F)\) est héréditaire pour tous les
constructeurs du langage. Elle est donc vraie pour toute formule
\(F \in \mathcal{F}\).

\begin{center}\rule{0.5\linewidth}{0.5pt}\end{center}

\hypertarget{question-2-uxe9quivalence-de-validituxe9}{%
\subsubsection{Question 2 : Équivalence de
validité}\label{question-2-uxe9quivalence-de-validituxe9}}

Soit \(F\) une formule contenant les variables \(p_1, \dots, p_n\). Le
résultat de la question 1 montre que la table de vérité de \(F^*\)
s'obtient en prenant le complément de la table de vérité de \(F\)
évaluée sur les lignes conjuguées. Si l'on substitue uniformément chaque
variable \(p_i\) par sa négation dans \(F^*\), on compense le passage à
la valuation conjuguée \(\bar{v}\). Ainsi, en posant
\(H = \neg (F^*[\neg p_1/p_1, \dots, \neg p_n/p_n])\), nous avons pour
toute valuation \(v\) :
\[v(H) = 1 - v(F^*[\neg p_1/p_1, \dots, \neg p_n/p_n]) = 1 - (1 - v(F)) = v(F)\]
Puisque \(v(H) = v(F)\) pour toute valuation \(v\), la formule \(H\) est
sémantiquement équivalente à \(F\). Par conséquent, \(F\) est une
tautologie si et seulement si \(H\) est une tautologie.

\begin{center}\rule{0.5\linewidth}{0.5pt}\end{center}

\hypertarget{question-3-duale-de-la-formule-g-p-land-q-lor-r}{%
\subsubsection{\texorpdfstring{Question 3 : Duale de la formule
\(G = (P \land Q) \lor R\)}{Question 3 : Duale de la formule G = (P \textbackslash land Q) \textbackslash lor R}}\label{question-3-duale-de-la-formule-g-p-land-q-lor-r}}

Appliquons pas-à-pas les règles de construction récursive du dual : 1.
Repérons le connecteur principal de \(G\). Il s'agit d'une disjonction
\(\lor\) entre la sous-formule \(A = P \land Q\) et \(B = R\). 2. Par la
règle 4 : \[G^* = (P \land Q)^* \land R^*\] 3. Déterminons
\((P \land Q)^*\) par la règle 3 : \[(P \land Q)^* = P^* \lor Q^*\] 4.
D'après la règle 1 pour les atomes :
\[P^* = P, \quad Q^* = Q, \quad R^* = R\] 5. Réassemblons le tout :
\[G^* = (P \lor Q) \land R\]

La duale de \((P \land Q) \lor R\) est donc \((P \lor Q) \land R\). On
remarque que la conjonction et la disjonction ont été permutées.
\n
\hypertarget{exercice-4-preuve-en-duxe9duction-naturelle-de-la-loi-de-peirce}{%
\section{Exercice 4 : Preuve en déduction naturelle de la loi de
Peirce}\label{exercice-4-preuve-en-duxe9duction-naturelle-de-la-loi-de-peirce}}

\hypertarget{uxe9noncuxe9}{%
\subsection{Énoncé}\label{uxe9noncuxe9}}

La \textbf{déduction naturelle} est un système formel syntaxique
introduit par Gerhard Gentzen. Dans sa version classique, il dispose
notamment des règles suivantes (où \(\Gamma\) est un ensemble de
formules servant d'hypothèses) : - \textbf{Règle
d'Implication-Introduction (\(\Rightarrow\)-I) :} Si
\(\Gamma, A \vdash B\), alors \(\Gamma \vdash A \Rightarrow B\). -
\textbf{Règle d'Implication-Élimination ou Modus Ponens
(\(\Rightarrow\)-E) :} Si \(\Gamma \vdash A \Rightarrow B\) et
\(\Gamma \vdash A\), alors \(\Gamma \vdash B\). - \textbf{Règle de
l'Absurde Classique ou Élimination de la Double Négation
(\(\neg\neg\)-E) :} Si \(\Gamma, \neg A \vdash \bot\) (où \(\bot\)
représente une contradiction), alors \(\Gamma \vdash A\). -
\textbf{Règle de Négation-Introduction (\(\neg\)-I) :} Si
\(\Gamma, A \vdash \bot\), alors \(\Gamma \vdash \neg A\). -
\textbf{Règle de Négation-Élimination (\(\neg\)-E) :} Si
\(\Gamma \vdash \neg A\) et \(\Gamma \vdash A\), alors
\(\Gamma \vdash \bot\).

Démontrer la \textbf{loi de Peirce}
\(\vdash ((A \Rightarrow B) \Rightarrow A) \Rightarrow A\) dans ce
système formel en détaillant chaque étape de déduction.

\begin{center}\rule{0.5\linewidth}{0.5pt}\end{center}

\hypertarget{correction-duxe9tailluxe9e}{%
\subsection{Correction Détaillée}\label{correction-duxe9tailluxe9e}}

Pour démontrer la formule
\(((A \Rightarrow B) \Rightarrow A) \Rightarrow A\), nous allons
chercher à appliquer la règle d'introduction de l'implication
(\(\Rightarrow\)-I) à l'étape finale. Il nous suffit donc de prouver le
séquent suivant sous hypothèse :
\[(A \Rightarrow B) \Rightarrow A \vdash A\]

Pour démontrer \(A\) sous cette hypothèse, nous allons procéder par
l'absurde classique. Nous rajoutons donc l'hypothèse \(\neg A\) et
cherchons à déduire la contradiction \(\bot\). Notre but intermédiaire
est de prouver :
\[(A \Rightarrow B) \Rightarrow A, \quad \neg A \vdash \bot\]

Voici la construction pas-à-pas de la démonstration :

\hypertarget{uxe9tape-1-hypothuxe8ses-de-duxe9part}{%
\subsubsection{Étape 1 : Hypothèses de
départ}\label{uxe9tape-1-hypothuxe8ses-de-duxe9part}}

Nous disposons dans notre contexte de deux hypothèses : - \((1)\) :
\((A \Rightarrow B) \Rightarrow A\) - \((2)\) : \(\neg A\)

\hypertarget{uxe9tape-2-construction-de-la-sous-preuve-de-a-rightarrow-b}{%
\subsubsection{\texorpdfstring{Étape 2 : Construction de la sous-preuve
de
\(A \Rightarrow B\)}{Étape 2 : Construction de la sous-preuve de A \textbackslash Rightarrow B}}\label{uxe9tape-2-construction-de-la-sous-preuve-de-a-rightarrow-b}}

Pour exploiter l'hypothèse \((1)\) par élimination de l'implication, il
serait utile de disposer de la formule \(A \Rightarrow B\). Montrons que
sous les hypothèses \((1)\) et \((2)\), nous pouvons déduire
\(A \Rightarrow B\). Pour cela, introduisons une hypothèse temporaire
supplémentaire pour construire l'implication : - \((3)\) : \(A\)
(hypothèse pour l'introduction de \(\Rightarrow\))

Maintenant, nous avons dans le contexte local les formules \(A\)
(hypothèse 3) et \(\neg A\) (hypothèse 2). 1. Appliquons la règle de
négation-élimination (\(\neg\)-E) sur \(A\) et \(\neg A\) :
\[\text{De } \Gamma, A \vdash A \text{ et } \Gamma, A \vdash \neg A \implies \Gamma, A \vdash \bot\]
Nous obtenons une contradiction \(\bot\). 2. Sous la contradiction
\(\bot\), nous pouvons déduire n'importe quelle proposition (règle du
\emph{Principe d'explosion} ou \emph{ex falso sequitur quodlibet}).
Appliquons cette règle pour déduire \(B\) :
\[\text{Puisque } (A \Rightarrow B) \Rightarrow A, \neg A, A \vdash \bot \implies (A \Rightarrow B) \Rightarrow A, \neg A, A \vdash B\]
3. Déchargeons maintenant l'hypothèse temporaire \(A\) (3) par la règle
d'introduction de l'implication (\(\Rightarrow\)-I) :
\[\text{Puisque } (A \Rightarrow B) \Rightarrow A, \neg A, A \vdash B \implies (A \Rightarrow B) \Rightarrow A, \neg A \vdash A \Rightarrow B\]
Nous avons réussi à démontrer la formule \(A \Rightarrow B\) sous les
hypothèses de base.

\hypertarget{uxe9tape-3-application-du-modus-ponens-et-obtention-de-la-contradiction}{%
\subsubsection{Étape 3 : Application du Modus Ponens et Obtention de la
Contradiction}\label{uxe9tape-3-application-du-modus-ponens-et-obtention-de-la-contradiction}}

Reprenons notre contexte global avec les hypothèses \((1)\)
\((A \Rightarrow B) \Rightarrow A\) et \((2)\) \(\neg A\). 1. Nous
venons de prouver que ce contexte permet de déduire \(A \Rightarrow B\).
2. Appliquons la règle d'élimination de l'implication
(\(\Rightarrow\)-E, ou Modus Ponens) entre l'hypothèse \((1)\) et la
formule \(A \Rightarrow B\) :
\[\text{De } \Gamma \vdash (A \Rightarrow B) \Rightarrow A \quad \text{et} \quad \Gamma \vdash A \Rightarrow B \implies \Gamma \vdash A\]
Nous avons donc déduit \(A\). 3. Nous disposons désormais à la fois de
\(A\) et de l'hypothèse de base \(\neg A\) (2). Appliquons la règle de
négation-élimination (\(\neg\)-E) :
\[\text{De } \Gamma \vdash A \quad \text{et} \quad \Gamma \vdash \neg A \implies \Gamma \vdash \bot\]
Nous obtenons la contradiction souhaitée.

\hypertarget{uxe9tape-4-duxe9chargement-des-hypothuxe8ses}{%
\subsubsection{Étape 4 : Déchargement des
hypothèses}\label{uxe9tape-4-duxe9chargement-des-hypothuxe8ses}}

\begin{enumerate}
\def\labelenumi{\arabic{enumi}.}
\tightlist
\item
  Appliquons la règle de l'absurde classique (\(\neg\neg\)-E) pour
  décharger l'hypothèse \(\neg A\) (2) :
  \[\text{Puisque } (A \Rightarrow B) \Rightarrow A, \neg A \vdash \bot \implies (A \Rightarrow B) \Rightarrow A \vdash A\]
\item
  Appliquons enfin la règle d'introduction de l'implication
  (\(\Rightarrow\)-I) pour décharger la dernière hypothèse
  \((A \Rightarrow B) \Rightarrow A\) (1) :
  \[\text{Puisque } (A \Rightarrow B) \Rightarrow A \vdash A \implies \vdash ((A \Rightarrow B) \Rightarrow A) \Rightarrow A\]
\end{enumerate}

La loi de Peirce est démontrée de manière purement syntaxique dans le
système classique de déduction naturelle.

\begin{center}\rule{0.5\linewidth}{0.5pt}\end{center}

\hypertarget{repruxe9sentation-sous-forme-darbre-de-preuve}{%
\subsubsection{Représentation sous forme d'Arbre de
Preuve}\label{repruxe9sentation-sous-forme-darbre-de-preuve}}

Voici l'arbre syntaxique complet résumant la démonstration :

\[\frac{
  \frac{
    \displaystyle [(A \Rightarrow B) \Rightarrow A]^1 \qquad
    \frac{
      \frac{\displaystyle [\neg A]^2 \quad [A]^3}{\bot}\small\text{($\neg$-E)}
      }{
        \displaystyle B
      }\small\text{($\bot$-E)}
    }{
      \displaystyle A \Rightarrow B
    }\small\text{($\Rightarrow$-I décharge 3)}
  }{
    \displaystyle A
  }\small\text{($\Rightarrow$-E)} \qquad [\neg A]^2
}{
  \frac{
    \displaystyle \bot
  }{
    \displaystyle A
  }\small\text{($\neg\neg$-E décharge 2)}
}\small\text{($\neg$-E)}\]
\[\overline{\displaystyle ((A \Rightarrow B) \Rightarrow A) \Rightarrow A}\small\text{($\Rightarrow$-I décharge 1)}\]
\n
\hypertarget{exercice-5-conversion-en-forme-normale-conjonctive-cnf}{%
\section{Exercice 5 : Conversion en forme normale conjonctive
(CNF)}\label{exercice-5-conversion-en-forme-normale-conjonctive-cnf}}

\hypertarget{uxe9noncuxe9}{%
\subsection{Énoncé}\label{uxe9noncuxe9}}

Soient \(P\), \(Q\), \(R\) et \(S\) quatre variables propositionnelles.
On considère la formule propositionnelle suivante :
\[H = (P \Rightarrow Q) \Rightarrow (R \land \neg S)\]

Convertir la formule \(H\) en \textbf{Forme Normale Conjonctive (CNF)}
(conjonction de clauses disjonctives) en détaillant chaque étape de
transformation algébrique.

\begin{center}\rule{0.5\linewidth}{0.5pt}\end{center}

\hypertarget{correction-duxe9tailluxe9e}{%
\subsection{Correction Détaillée}\label{correction-duxe9tailluxe9e}}

Une formule est sous Forme Normale Conjonctive (CNF) si elle se présente
sous la forme :
\[\bigwedge_{i=1}^n \left( \bigvee_{j=1}^{m_i} L_{i,j} \right)\] où
chaque \(L_{i,j}\) est un littéral (une variable propositionnelle ou sa
négation).

Pour transformer \(H\) en CNF, nous appliquons un algorithme en trois
grandes étapes : 1. Élimination des connecteurs d'implication
(\(\Rightarrow\)). 2. Descente des négations (\(\neg\)) au niveau des
variables (en utilisant les lois de De Morgan et de double négation). 3.
Distribution des disjonctions (\(\lor\)) par rapport aux conjonctions
(\(\land\)).

\begin{center}\rule{0.5\linewidth}{0.5pt}\end{center}

\hypertarget{uxe9tape-1-uxe9limination-des-connecteurs-dimplication-rightarrow}{%
\subsubsection{\texorpdfstring{Étape 1 : Élimination des connecteurs
d'implication
(\(\Rightarrow\))}{Étape 1 : Élimination des connecteurs d'implication (\textbackslash Rightarrow)}}\label{uxe9tape-1-uxe9limination-des-connecteurs-dimplication-rightarrow}}

Rappelons la règle d'équivalence de base :
\(A \Rightarrow B \equiv \neg A \lor B\).

\begin{enumerate}
\def\labelenumi{\arabic{enumi}.}
\tightlist
\item
  Éliminons l'implication principale de \(H\). Posons
  \(A = (P \Rightarrow Q)\) et \(B = (R \land \neg S)\). L'expression
  devient : \[H \equiv \neg(P \Rightarrow Q) \lor (R \land \neg S)\]
\item
  Éliminons l'implication résiduelle à l'intérieur de la négation :
  \[P \Rightarrow Q \equiv \neg P \lor Q\]
\item
  Substituons ce résultat dans la formule complète :
  \[H \equiv \neg(\neg P \lor Q) \lor (R \land \neg S)\]
\end{enumerate}

\begin{center}\rule{0.5\linewidth}{0.5pt}\end{center}

\hypertarget{uxe9tape-2-descente-des-nuxe9gations-neg}{%
\subsubsection{\texorpdfstring{Étape 2 : Descente des négations
(\(\neg\))}{Étape 2 : Descente des négations (\textbackslash neg)}}\label{uxe9tape-2-descente-des-nuxe9gations-neg}}

Appliquons la loi de De Morgan sur le premier bloc
\(\neg(\neg P \lor Q)\) :
\[\neg(\neg P \lor Q) \equiv \neg(\neg P) \land \neg Q\]

Par la règle de double négation (\(\neg\neg P \equiv P\)) :
\[\neg(\neg P \lor Q) \equiv P \land \neg Q\]

Substituons cette forme simplifiée dans l'expression globale de \(H\) :
\[H \equiv (P \land \neg Q) \lor (R \land \neg S)\]

\begin{center}\rule{0.5\linewidth}{0.5pt}\end{center}

\hypertarget{uxe9tape-3-distribution-de-la-disjonction-lor-par-rapport-uxe0-la-conjonction-land}{%
\subsubsection{\texorpdfstring{Étape 3 : Distribution de la disjonction
(\(\lor\)) par rapport à la conjonction
(\(\land\))}{Étape 3 : Distribution de la disjonction (\textbackslash lor) par rapport à la conjonction (\textbackslash land)}}\label{uxe9tape-3-distribution-de-la-disjonction-lor-par-rapport-uxe0-la-conjonction-land}}

Nous devons maintenant distribuer la disjonction extérieure par rapport
aux conjonctions intérieures. Soient \(X = (P \land \neg Q)\) et
\(Y = (R \land \neg S)\). L'expression est \(X \lor Y\).

\begin{enumerate}
\def\labelenumi{\arabic{enumi}.}
\tightlist
\item
  Distribuons d'abord le bloc \(X\) sur la conjonction
  \(Y = (R \land \neg S)\) :
  \[X \lor (R \land \neg S) \equiv (X \lor R) \land (X \lor \neg S)\]
\item
  Remplaçons maintenant \(X\) par sa valeur originelle
  \((P \land \neg Q)\) dans le premier terme \((X \lor R)\) :
  \[X \lor R \equiv (P \land \neg Q) \lor R\] Appliquons la
  distributivité de \(\lor\) par rapport à \(\land\) sur ce terme :
  \[(P \land \neg Q) \lor R \equiv (P \lor R) \land (\neg Q \lor R)\]
\item
  Remplaçons \(X\) dans le second terme \((X \lor \neg S)\) :
  \[X \lor \neg S \equiv (P \land \neg Q) \lor \neg S\] Distribuons
  également :
  \[(P \land \neg Q) \lor \neg S \equiv (P \lor \neg S) \land (\neg Q \lor \neg S)\]
\item
  Assemblons maintenant les résultats des points 2 et 3 dans la formule
  générale obtenue au point 1 :
  \[H \equiv \Big( (P \lor R) \land (\neg Q \lor R) \Big) \land \Big( (P \lor \neg S) \land (\neg Q \lor \neg S) \Big)\]
\item
  Par associativité de la conjonction \(\land\), nous pouvons supprimer
  les parenthèses superflues :
  \[H \equiv (P \lor R) \land (\neg Q \lor R) \land (P \lor \neg S) \land (\neg Q \lor \neg S)\]
\end{enumerate}

\begin{center}\rule{0.5\linewidth}{0.5pt}\end{center}

\hypertarget{ruxe9sultat-final}{%
\subsubsection{Résultat final}\label{ruxe9sultat-final}}

La formule \(H\) sous Forme Normale Conjonctive est :
\[H \equiv (P \lor R) \land (\neg Q \lor R) \land (P \lor \neg S) \land (\neg Q \lor \neg S)\]

Cette CNF est composée de \(4\) clauses disjonctives : - Clause 1 :
\(P \lor R\) - Clause 2 : \(\neg Q \lor R\) - Clause 3 :
\(P \lor \neg S\) - Clause 4 : \(\neg Q \lor \neg S\)
\n
\hypertarget{exercice-6-ruxe9futation-et-ruxe9solution-de-clauses}{%
\section{Exercice 6 : Réfutation et résolution de
clauses}\label{exercice-6-ruxe9futation-et-ruxe9solution-de-clauses}}

\hypertarget{uxe9noncuxe9}{%
\subsection{Énoncé}\label{uxe9noncuxe9}}

En logique propositionnelle, la \textbf{résolution} est une règle
d'inférence syntaxique qui opère sur des clauses. La règle de résolution
s'énonce ainsi :
\[\text{De } A \lor B \quad \text{et} \quad \neg A \lor C, \quad \text{on déduit} \quad B \lor C\]
La clause \(B \lor C\) est appelée la \textbf{résolvante} des deux
clauses parentes. La variable \(A\) est la variable pivot. Si l'une des
clauses est réduite à un littéral unique (par exemple \(A\)) et l'autre
à sa négation (par exemple \(\neg A\)), la résolvante est la
\textbf{clause vide}, notée \(\square\) (ou \(\bot\)), qui représente
une contradiction.

Considérons la base de connaissances \(\Sigma\) composée des quatre
clauses suivantes : 1. \(C_1 = P \lor Q\) 2. \(C_2 = \neg P \lor Q\) 3.
\(C_3 = P \lor \neg Q\) 4. \(C_4 = \neg P \lor \neg Q\)

\begin{enumerate}
\def\labelenumi{\arabic{enumi}.}
\tightlist
\item
  Expliquer le principe de la \textbf{démonstration par réfutation}
  (preuve par l'absurde à l'aide de la résolution).
\item
  Utiliser la méthode de réfutation par résolution pour démontrer que la
  base de connaissances \(\Sigma\) est incompatible (c'est-à-dire non
  satisfaisable). On construira pour cela un arbre de résolution menant
  à la clause vide \(\square\).
\end{enumerate}

\begin{center}\rule{0.5\linewidth}{0.5pt}\end{center}

\hypertarget{correction-duxe9tailluxe9e}{%
\subsection{Correction Détaillée}\label{correction-duxe9tailluxe9e}}

\hypertarget{question-1-principe-de-la-duxe9monstration-par-ruxe9futation}{%
\subsubsection{Question 1 : Principe de la démonstration par
réfutation}\label{question-1-principe-de-la-duxe9monstration-par-ruxe9futation}}

La démonstration par réfutation s'appuie sur le théorème logique suivant
: un ensemble de formules \(\Sigma\) implique logiquement une formule
\(A\) (noté \(\Sigma \models A\)) si et seulement si l'ensemble
\(\Sigma \cup \{\neg A\}\) est insatisfaisable (c'est-à-dire
contradictoire).

En résolution, pour démontrer qu'une base de connaissances est
insatisfaisable, on applique de manière répétée la règle d'inférence de
résolution sur les clauses disponibles dans \(\Sigma\). Si l'on parvient
à déduire la clause vide \(\square\), cela signifie que la base de
connaissances contient une contradiction sémantique sous-jacente. Comme
la règle de résolution est \textbf{correcte} (si elle produit
\(\square\), alors l'ensemble est insatisfaisable) et \textbf{complète
pour la réfutation} (si l'ensemble est insatisfaisable, il existe
toujours une suite d'applications de la résolution menant à
\(\square\)), l'obtention de la clause vide constitue une preuve
irréfutable d'insatisfaisabilité.

\begin{center}\rule{0.5\linewidth}{0.5pt}\end{center}

\hypertarget{question-2-ruxe9futation-par-ruxe9solution-sur-sigma}{%
\subsubsection{\texorpdfstring{Question 2 : Réfutation par résolution
sur
\(\Sigma\)}{Question 2 : Réfutation par résolution sur \textbackslash Sigma}}\label{question-2-ruxe9futation-par-ruxe9solution-sur-sigma}}

Notre but est de déduire la clause vide \(\square\) à partir de notre
ensemble de clauses de départ :
\[\Sigma = \{ P \lor Q, \quad \neg P \lor Q, \quad P \lor \neg Q, \quad \neg P \lor \neg Q \}\]

Appliquons la règle de résolution pas-à-pas en choisissant
judicieusement les pivots :

\begin{enumerate}
\def\labelenumi{\arabic{enumi}.}
\item
  \textbf{Étape 1 : Résolution sur le pivot \(P\) (première paire de
  clauses)} Considérons les clauses parentes :

  \begin{itemize}
  \tightlist
  \item
    \(C_1 = P \lor Q\)
  \item
    \(C_2 = \neg P \lor Q\)
  \end{itemize}

  La variable pivot est \(P\). La règle de résolution élimine \(P\) et
  \(\neg P\) pour former la disjonction des littéraux restants :
  \[C_5 = \text{Res}(C_1, C_2) = Q \lor Q\] Par la règle d'idempotence
  de la disjonction (\(Q \lor Q \equiv Q\)), nous obtenons le littéral
  unitaire : \[C_5 = Q\]
\item
  \textbf{Étape 2 : Résolution sur le pivot \(P\) (seconde paire de
  clauses)} Considérons les clauses parentes :

  \begin{itemize}
  \tightlist
  \item
    \(C_3 = P \lor \neg Q\)
  \item
    \(C_4 = \neg P \lor \neg Q\)
  \end{itemize}

  La variable pivot est \(P\). La résolution produit :
  \[C_6 = \text{Res}(C_3, C_4) = \neg Q \lor \neg Q\] Par la règle
  d'idempotence, nous obtenons le littéral unitaire : \[C_6 = \neg Q\]
\item
  \textbf{Étape 3 : Déduction de la contradiction (clause vide)}
  Considérons désormais les deux clauses dérivées unitaires :

  \begin{itemize}
  \tightlist
  \item
    \(C_5 = Q\)
  \item
    \(C_6 = \neg Q\)
  \end{itemize}

  Appliquons la règle de résolution sur le pivot \(Q\). Puisqu'il n'y a
  plus aucun autre littéral dans les clauses, nous déduisons la clause
  vide : \[C_7 = \text{Res}(C_5, C_6) = \square\]
\end{enumerate}

L'arbre de résolution menant à la clause vide \(\square\) est donc
construit avec succès. La base de connaissances \(\Sigma\) est
incompatible et non satisfaisable.

\begin{center}\rule{0.5\linewidth}{0.5pt}\end{center}

\hypertarget{repruxe9sentation-graphique-de-larbre-de-ruxe9solution}{%
\subsubsection{Représentation graphique de l'arbre de
résolution}\label{repruxe9sentation-graphique-de-larbre-de-ruxe9solution}}

\begin{lstlisting}
\NormalTok{       C1: P v Q     C2: \textasciitilde{}P v Q           C3: P v \textasciitilde{}Q    C4: \textasciitilde{}P v \textasciitilde{}Q}
\NormalTok{          \textbackslash{}           /                      \textbackslash{}           /}
\NormalTok{           \textbackslash{}         /                        \textbackslash{}         /}
\NormalTok{            \textbackslash{}       /                          \textbackslash{}       /}
\NormalTok{             C5: Q                              C6: \textasciitilde{}Q}
\NormalTok{               \textbackslash{}                                  /}
\NormalTok{                \textbackslash{}                                /}
\NormalTok{                 \textbackslash{}                              /}
\NormalTok{                  \textbackslash{}                            /}
\NormalTok{                         C7: [Clause vide]}
\end{lstlisting}
\n
\hypertarget{exercice-7-compluxe9tude-du-connecteur-nor-barre-de-peirce}{%
\section{Exercice 7 : Complétude du connecteur NOR (barre de
Peirce)}\label{exercice-7-compluxe9tude-du-connecteur-nor-barre-de-peirce}}

\hypertarget{uxe9noncuxe9}{%
\subsection{Énoncé}\label{uxe9noncuxe9}}

Le connecteur binaire \(\downarrow\) (appelé connecteur NOR ou flèche de
Peirce) est défini par la règle d'évaluation sémantique suivante :
\[\text{Pour toute valuation } v, \quad v(A \downarrow B) = 1 \iff v(A) = 0 \text{ et } v(B) = 0\]
Ce qui correspond sémantiquement à \(\neg(A \lor B)\).

Démontrer que le singleton \(\{\downarrow\}\) est un \textbf{système
complet de connecteurs} (ou base fonctionnellement complète) en
exprimant à l'aide de ce seul connecteur : 1. La négation \(\neg A\). 2.
La disjonction \(A \lor B\). 3. La conjonction \(A \land B\).

\begin{center}\rule{0.5\linewidth}{0.5pt}\end{center}

\hypertarget{correction-duxe9tailluxe9e}{%
\subsection{Correction Détaillée}\label{correction-duxe9tailluxe9e}}

\hypertarget{question-1-expression-de-la-nuxe9gation-neg-a}{%
\subsubsection{\texorpdfstring{Question 1 : Expression de la négation
\(\neg A\)}{Question 1 : Expression de la négation \textbackslash neg A}}\label{question-1-expression-de-la-nuxe9gation-neg-a}}

Montrons que \(\neg A \equiv A \downarrow A\). Soit \(v\) une valuation
quelconque. - D'après la définition du connecteur \(\downarrow\) :
\[v(A \downarrow A) = 1 \iff v(A) = 0 \text{ et } v(A) = 0 \iff v(A) = 0\]
- Ainsi, \(v(A \downarrow A) = 1 - v(A) = v(\neg A)\). Les deux
expressions prennent les mêmes valeurs de vérité pour toutes les
valuations possibles. Nous avons donc bien :
\[\neg A \equiv A \downarrow A\]

\begin{center}\rule{0.5\linewidth}{0.5pt}\end{center}

\hypertarget{question-2-expression-de-la-disjonction-a-lor-b}{%
\subsubsection{\texorpdfstring{Question 2 : Expression de la disjonction
\(A \lor B\)}{Question 2 : Expression de la disjonction A \textbackslash lor B}}\label{question-2-expression-de-la-disjonction-a-lor-b}}

Par double négation, nous savons que
\(A \lor B \equiv \neg(\neg(A \lor B))\). - Par définition du connecteur
\(\downarrow\), nous avons \(\neg(A \lor B) \equiv A \downarrow B\). -
Par conséquent, la disjonction est la négation du connecteur
\(\downarrow\) : \[A \lor B \equiv \neg(A \downarrow B)\] - En utilisant
le résultat de la question 1 pour exprimer la négation d'une formule
\(X = (A \downarrow B)\), nous avons :
\[\neg X \equiv X \downarrow X \implies \neg(A \downarrow B) \equiv (A \downarrow B) \downarrow (A \downarrow B)\]
Les deux expressions sont logiquement équivalentes :
\[A \lor B \equiv (A \downarrow B) \downarrow (A \downarrow B)\]

\begin{center}\rule{0.5\linewidth}{0.5pt}\end{center}

\hypertarget{question-3-expression-de-la-conjonction-a-land-b}{%
\subsubsection{\texorpdfstring{Question 3 : Expression de la conjonction
\(A \land B\)}{Question 3 : Expression de la conjonction A \textbackslash land B}}\label{question-3-expression-de-la-conjonction-a-land-b}}

Par les lois de De Morgan, nous savons que la conjonction peut s'écrire
sous forme d'une négation sur des termes niés :
\[A \land B \equiv \neg(\neg A \lor \neg B)\] - Par définition du
connecteur \(\downarrow\), la négation d'une disjonction s'écrit avec la
flèche : \[\neg(X \lor Y) \equiv X \downarrow Y\] - Posons
\(X = \neg A\) et \(Y = \neg B\). L'expression devient :
\[A \land B \equiv \neg A \downarrow \neg B\] - Remplaçons désormais les
négations internes \(\neg A\) et \(\neg B\) en utilisant le connecteur
\(\downarrow\) d'après la question 1 :
\[\neg A \equiv A \downarrow A \quad \text{et} \quad \neg B \equiv B \downarrow B\]
- En substituant ces formes, nous obtenons :
\[A \land B \equiv (A \downarrow A) \downarrow (B \downarrow B)\]

\begin{center}\rule{0.5\linewidth}{0.5pt}\end{center}

\hypertarget{synthuxe8se}{%
\subsubsection{Synthèse}\label{synthuxe8se}}

Nous avons réussi à exprimer les trois connecteurs fondamentaux
\(\{\neg, \lor, \land\}\) uniquement à l'aide de la flèche de Peirce
\(\downarrow\) : - \(\neg A \equiv A \downarrow A\) -
\(A \lor B \equiv (A \downarrow B) \downarrow (A \downarrow B)\) -
\(A \land B \equiv (A \downarrow A) \downarrow (B \downarrow B)\)

Puisque le système de connecteurs \(\{\neg, \lor, \land\}\) est complet,
le singleton \(\{\downarrow\}\) est également un système complet de
connecteurs.
\n
\hypertarget{exercice-8-thuxe9oruxe8me-dinterpolation-de-craig}{%
\section{Exercice 8 : Théorème d'interpolation de
Craig}\label{exercice-8-thuxe9oruxe8me-dinterpolation-de-craig}}

\hypertarget{uxe9noncuxe9}{%
\subsection{Énoncé}\label{uxe9noncuxe9}}

Le \textbf{théorème d'interpolation de Craig} est un résultat
fondamental de la logique mathématique. Dans le cadre de la logique
propositionnelle, il s'énonce ainsi : \textgreater{} Soient \(A\) et
\(B\) deux formules telles que \(A \Rightarrow B\) est une tautologie
(\(\models A \Rightarrow B\)). \textgreater{} Si \(A\) et \(B\)
partagent au moins une variable propositionnelle commune, alors il
existe une formule \(I\), appelée \textbf{interpolant de Craig}, telle
que : \textgreater{} 1. \(\models A \Rightarrow I\) \textgreater{} 2.
\(\models I \Rightarrow B\) \textgreater{} 3. Toutes les variables
propositionnelles apparaissant dans \(I\) apparaissent à la fois dans
\(A\) et dans \(B\).

Considérons les deux formules propositionnelles suivantes : -
\(A = P \land (\neg Q \lor R)\) -
\(B = (P \land S) \lor (P \land \neg S \lor T)\)

\begin{enumerate}
\def\labelenumi{\arabic{enumi}.}
\tightlist
\item
  Démontrer que \(\models A \Rightarrow B\) en effectuant une
  simplification sémantique de \(A\) et \(B\).
\item
  Déterminer l'ensemble \(\mathcal{V}_A\) des variables de \(A\),
  l'ensemble \(\mathcal{V}_B\) des variables de \(B\), et leur
  intersection \(\mathcal{V}_C = \mathcal{V}_A \cap \mathcal{V}_B\).
\item
  Proposer un interpolant de Craig \(I\) pour ce couple de formules et
  prouver qu'il satisfait aux trois conditions du théorème.
\end{enumerate}

\begin{center}\rule{0.5\linewidth}{0.5pt}\end{center}

\hypertarget{correction-duxe9tailluxe9e}{%
\subsection{Correction Détaillée}\label{correction-duxe9tailluxe9e}}

\hypertarget{question-1-preuve-de-la-tautologie-models-a-rightarrow-b}{%
\subsubsection{\texorpdfstring{Question 1 : Preuve de la tautologie
\(\models A \Rightarrow B\)}{Question 1 : Preuve de la tautologie \textbackslash models A \textbackslash Rightarrow B}}\label{question-1-preuve-de-la-tautologie-models-a-rightarrow-b}}

Commençons par simplifier et analyser les conditions sous lesquelles les
deux formules sont vraies.

\begin{itemize}
\item
  \textbf{Analyse de \(A\) :} \[A = P \land (\neg Q \lor R)\] Pour que
  \(v(A) = 1\), il faut et il suffit que :

  \begin{itemize}
  \tightlist
  \item
    \(v(P) = 1\)
  \item
    et \(v(\neg Q \lor R) = 1\). En particulier, toute valuation
    satisfaisant \(A\) doit nécessairement affecter la valeur \(1\) à la
    variable \(P\).
  \end{itemize}
\item
  \textbf{Analyse de \(B\) :}
  \[B = (P \land S) \lor (P \land \neg S \lor T)\] Grâce à la
  commutativité et l'associativité de la disjonction, réécrivons \(B\) :
  \[B \equiv (P \land S) \lor (P \land \neg S) \lor T\] Appliquons la
  distributivité de la conjonction par rapport à la disjonction sur les
  deux premiers termes :
  \[(P \land S) \lor (P \land \neg S) \equiv P \land (S \lor \neg S)\]
  Puisque \(S \lor \neg S \equiv 1\), nous avons :
  \[P \land 1 \equiv P\] Substituons cela dans l'expression de \(B\) :
  \[B \equiv P \lor T\]
\item
  \textbf{Analyse de l'implication \(A \Rightarrow B\) :} Remplaçons
  \(B\) par sa forme simplifiée \(P \lor T\) :
  \[A \Rightarrow B \equiv \big( P \land (\neg Q \lor R) \big) \Rightarrow (P \lor T)\]
  Soit \(v\) une valuation telle que \(v(A) = 1\). D'après l'analyse de
  \(A\), cela implique \(v(P) = 1\). Si \(v(P) = 1\), alors la
  disjonction \(P \lor T\) est vraie :
  \[v(P \lor T) = \max(v(P), v(T)) = \max(1, v(T)) = 1 \implies v(B) = 1\]
  Ainsi, toute valuation qui rend \(A\) vraie rend également \(B\)
  vraie. Par conséquent, \(A \Rightarrow B\) est une tautologie, notée
  \(\models A \Rightarrow B\).
\end{itemize}

\begin{center}\rule{0.5\linewidth}{0.5pt}\end{center}

\hypertarget{question-2-intersection-des-variables-propositionnelles}{%
\subsubsection{Question 2 : Intersection des variables
propositionnelles}\label{question-2-intersection-des-variables-propositionnelles}}

Listons les variables propositionnelles contenues dans chaque formule :
- Pour \(A = P \land (\neg Q \lor R)\), l'ensemble des variables est :
\[\mathcal{V}_A = \{P, Q, R\}\] - Pour
\(B = (P \land S) \lor (P \land \neg S \lor T)\), l'ensemble des
variables est : \[\mathcal{V}_B = \{P, S, T\}\] - L'intersection des
ensembles de variables est :
\[\mathcal{V}_C = \mathcal{V}_A \cap \mathcal{V}_B = \{P\}\]

L'unique variable commune aux deux formules est donc \(P\).
L'interpolant de Craig \(I\) ne pourra contenir que la seule variable
propositionnelle \(P\) (ou être une constante logique \(0\) ou \(1\)).

\begin{center}\rule{0.5\linewidth}{0.5pt}\end{center}

\hypertarget{question-3-duxe9termination-et-validation-de-linterpolant-i}{%
\subsubsection{\texorpdfstring{Question 3 : Détermination et validation
de l'interpolant
\(I\)}{Question 3 : Détermination et validation de l'interpolant I}}\label{question-3-duxe9termination-et-validation-de-linterpolant-i}}

Puisque le seul atome autorisé dans \(I\) est \(P\), les seuls choix
possibles pour \(I\) (à équivalence logique près) sont \(P\),
\(\neg P\), \(0\) ou \(1\).

\begin{itemize}
\tightlist
\item
  Testons la formule \textbf{\(I = P\)} :

  \begin{enumerate}
  \def\labelenumi{\arabic{enumi}.}
  \tightlist
  \item
    \textbf{Condition 3 :} Les variables de \(I\) sont incluses dans
    \(\mathcal{V}_C\). C'est le cas puisque
    \(\text{Var}(I) = \{P\} \subseteq \{P\}\).
  \item
    \textbf{Condition 1 (\(\models A \Rightarrow I\)) :}
    \[A \Rightarrow I \equiv \big( P \land (\neg Q \lor R) \big) \Rightarrow P\]
    Soit \(v\) une valuation. Si \(v(A) = 1\), alors \(v(P) = 1\), donc
    \(v(I) = 1\). L'implication est donc toujours vraie. C'est une
    tautologie.
  \item
    \textbf{Condition 2 (\(\models I \Rightarrow B\)) :}
    \[I \Rightarrow B \equiv P \Rightarrow (P \lor T)\] Soit \(v\) une
    valuation. Si \(v(I) = 1\), alors \(v(P) = 1\), ce qui implique
    \(v(P \lor T) = 1\), d'où \(v(B) = 1\). L'implication est donc
    toujours vraie. C'est une tautologie.
  \end{enumerate}
\end{itemize}

L'ensemble des trois conditions du théorème d'interpolation de Craig est
satisfait par la formule : \[I = P\]

\emph{Remarque constructive :} Pour construire l'interpolant
sémantiquement dans le cas général pour une variable commune \(C\), on
peut écrire \(I = A[C/1] \lor A[C/0]\) ou \(I = A[C/1] \land A[C/0]\)
selon le sens de l'implication.
\n
\hypertarget{exercice-9-isomorphisme-de-stone-pour-les-alguxe8bres-de-boole-finies}{%
\section{Exercice 9 : Isomorphisme de Stone pour les algèbres de Boole
finies}\label{exercice-9-isomorphisme-de-stone-pour-les-alguxe8bres-de-boole-finies}}

\hypertarget{uxe9noncuxe9}{%
\subsection{Énoncé}\label{uxe9noncuxe9}}

Soit \((B, \lor, \land, \neg, 0, 1)\) une \textbf{algèbre de Boole
finie} (c'est-à-dire que le cardinal de l'ensemble \(B\) est fini). On
munit \(B\) de sa relation d'ordre canonique \(\le\) définie par :
\[\forall x, y \in B, \quad x \le y \iff x \land y = x \iff x \lor y = y\]

On définit un \textbf{atom} de \(B\) comme un élément \(a \in B\) tel
que \(a \neq 0\) et pour tout \(x \in B\) :
\[\text{si } x \le a, \quad \text{alors } x = 0 \text{ ou } x = a\] On
note \(\mathcal{A}t(B)\) l'ensemble de tous les atomes de \(B\).

Considérons l'application suivante, notée \(\Phi\) :
\[\Phi : \begin{cases} B \to \mathcal{P}(\mathcal{A}t(B)) \\ x \mapsto \{a \in \mathcal{A}t(B) \mid a \le x\} \end{cases}\]

Le but de cet exercice est de démontrer le \textbf{théorème de
représentation de Stone} dans le cas fini, à savoir que \(\Phi\) est un
isomorphisme d'algèbres de Boole, prouvant que toute algèbre de Boole
finie est isomorphe à l'algèbre des parties d'un ensemble.

\begin{enumerate}
\def\labelenumi{\arabic{enumi}.}
\tightlist
\item
  Soit \(x \in B\). Démontrer que si \(x \neq 0\), alors il existe au
  moins un atome \(a \in \mathcal{A}t(B)\) tel que \(a \le x\).
\item
  Démontrer que pour tout \(x \in B\), \(x = \bigvee_{a \in \Phi(x)} a\)
  (tout élément est le suprémum des atomes qu'il contient).
\item
  Démontrer que \(\Phi\) est injective.
\item
  Démontrer que pour tous \(x, y \in B\) :

  \begin{itemize}
  \tightlist
  \item
    \(\Phi(x \land y) = \Phi(x) \cap \Phi(y)\)
  \item
    \(\Phi(x \lor y) = \Phi(x) \cup \Phi(y)\)
  \item
    \(\Phi(\neg x) = \mathcal{A}t(B) \setminus \Phi(x)\)
  \end{itemize}
\item
  Conclure sur la nature de \(\Phi\) et en déduire le cardinal d'une
  algèbre de Boole finie.
\end{enumerate}

\begin{center}\rule{0.5\linewidth}{0.5pt}\end{center}

\hypertarget{correction-duxe9tailluxe9e}{%
\subsection{Correction Détaillée}\label{correction-duxe9tailluxe9e}}

\hypertarget{question-1-existence-dun-atome-infuxe9rieur}{%
\subsubsection{Question 1 : Existence d'un atome
inférieur}\label{question-1-existence-dun-atome-infuxe9rieur}}

Soit \(x \in B\) tel que \(x \neq 0\). Considérons l'ensemble des
éléments non nuls inférieurs ou égaux à \(x\) :
\[E_x = \{y \in B \mid y \neq 0 \text{ et } y \le x\}\] - Cet ensemble
est non vide car \(x \in E_x\) (puisque \(x \neq 0\) et \(x \le x\)). -
Puisque l'ensemble \(B\) est fini, l'ensemble \(E_x\) est également
fini. - Tout ensemble ordonné fini et non vide possède au moins un
élément minimal. Soit \(a\) un élément minimal de \(E_x\).

Montrons que \(a\) est un atome de \(B\) : - Par définition de \(E_x\),
\(a \neq 0\). - Soit \(z \in B\) tel que \(z \le a\). Par transitivité
de la relation d'ordre, puisque \(a \le x\), on a \(z \le x\). - Si
\(z \neq 0\), alors \(z \in E_x\). Comme \(z \le a\) et que \(a\) est
minimal dans \(E_x\), on a nécessairement \(z = a\). - Ainsi, pour tout
\(z \le a\), \(z = 0\) ou \(z = a\). L'élément \(a\) est donc un atome
de \(B\). Puisque \(a \in E_x\), on a \(a \le x\). Il existe donc un
atome sous tout élément non nul.

\begin{center}\rule{0.5\linewidth}{0.5pt}\end{center}

\hypertarget{question-2-tout-uxe9luxe9ment-est-la-borne-supuxe9rieure-de-ses-atomes}{%
\subsubsection{Question 2 : Tout élément est la borne supérieure de ses
atomes}\label{question-2-tout-uxe9luxe9ment-est-la-borne-supuxe9rieure-de-ses-atomes}}

Soit \(x \in B\). Définissons \(x' = \bigvee_{a \in \Phi(x)} a\). - Par
définition de \(\Phi(x)\), pour tout \(a \in \Phi(x)\), nous avons
\(a \le x\). Par définition de la borne supérieure, nous en déduisons :
\[x' \le x\] - Montrons par l'absurde que \(x = x'\). Supposons que
\(x \neq x'\). Alors, l'élément \(d = x \land \neg x'\) est non nul
(sinon \(x \le x'\)). Puisque \(d \neq 0\), par le résultat de la
question 1, il existe un atome \(a_0 \in \mathcal{A}t(B)\) tel que
\(a_0 \le d\). - D'une part, \(a_0 \le d \le x\), donc
\(a_0 \in \Phi(x)\). Par conséquent, \(a_0 \le x'\) (par définition de
\(x'\)). - D'autre part, \(a_0 \le d \le \neg x'\), ce qui implique
\(a_0 \land x' \le \neg x' \land x' = 0\). Donc \(a_0 \land x' = 0\). -
Or, comme \(a_0 \le x'\), nous avons \(a_0 \land x' = a_0\). - Nous en
déduisons \(a_0 = 0\), ce qui contredit le fait que \(a_0\) est un atome
(donc non nul). L'hypothèse \(x \neq x'\) est donc absurde. On en
conclut que \(x = \bigvee_{a \in \Phi(x)} a\).

\begin{center}\rule{0.5\linewidth}{0.5pt}\end{center}

\hypertarget{question-3-injectivituxe9-de-phi}{%
\subsubsection{\texorpdfstring{Question 3 : Injectivité de
\(\Phi\)}{Question 3 : Injectivité de \textbackslash Phi}}\label{question-3-injectivituxe9-de-phi}}

Soient \(x, y \in B\) tels que \(\Phi(x) = \Phi(y)\). D'après le
résultat de la question 2 :
\[x = \bigvee_{a \in \Phi(x)} a \quad \text{et} \quad y = \bigvee_{a \in \Phi(y)} a\]
Puisque \(\Phi(x) = \Phi(y)\), les deux suprémums portent sur le même
ensemble d'atomes, d'où : \[x = y\] L'application \(\Phi\) est donc
injective.

\begin{center}\rule{0.5\linewidth}{0.5pt}\end{center}

\hypertarget{question-4-pruxe9servation-des-opuxe9rations-logiques}{%
\subsubsection{Question 4 : Préservation des opérations
logiques}\label{question-4-pruxe9servation-des-opuxe9rations-logiques}}

\begin{enumerate}
\def\labelenumi{\arabic{enumi}.}
\item
  \textbf{Intersection/Conjonction :
  \(\Phi(x \land y) = \Phi(x) \cap \Phi(y)\)}
  \[a \in \Phi(x \land y) \iff a \le x \land y \iff a \le x \text{ et } a \le y \iff a \in \Phi(x) \text{ et } a \in \Phi(y) \iff a \in \Phi(x) \cap \Phi(y)\]
\item
  \textbf{Union/Disjonction : \(\Phi(x \lor y) = \Phi(x) \cup \Phi(y)\)}

  \begin{itemize}
  \tightlist
  \item
    Si \(a \in \Phi(x) \cup \Phi(y)\), alors \(a \le x\) ou \(a \le y\).
    Dans les deux cas, \(a \le x \lor y\), donc
    \(a \in \Phi(x \lor y)\). D'où
    \(\Phi(x) \cup \Phi(y) \subseteq \Phi(x \lor y)\).
  \item
    Réciproquement, soit \(a \in \Phi(x \lor y)\), donc
    \(a \le x \lor y\). Alors
    \(a = a \land (x \lor y) = (a \land x) \lor (a \land y)\). Puisque
    \(a\) est un atome, les éléments \(a \land x\) et \(a \land y\) sont
    soit \(0\), soit \(a\). Si les deux valaient \(0\), on aurait
    \(a = 0 \lor 0 = 0\), ce qui est exclu. Donc, au moins l'un des deux
    vaut \(a\). Si \(a \land x = a\), alors
    \(a \le x \implies a \in \Phi(x)\). Si \(a \land y = a\), alors
    \(a \le y \implies a \in \Phi(y)\). Dans tous les cas,
    \(a \in \Phi(x) \cup \Phi(y)\). D'où l'égalité
    \(\Phi(x \lor y) = \Phi(x) \cup \Phi(y)\).
  \end{itemize}
\item
  \textbf{Complément :
  \(\Phi(\neg x) = \mathcal{A}t(B) \setminus \Phi(x)\)}

  \begin{itemize}
  \tightlist
  \item
    Soit \(a \in \Phi(\neg x)\), donc \(a \le \neg x\). Si
    \(a \in \Phi(x)\), on aurait également \(a \le x\), donc
    \(a \le x \land \neg x = 0 \implies a = 0\), absurde. Donc
    \(a \notin \Phi(x)\).
  \item
    Réciproquement, soit \(a \in \mathcal{A}t(B) \setminus \Phi(x)\).
    Considérons \(a \land x\). Puisque \(a\) est un atome,
    \(a \land x = 0\) ou \(a \land x = a\). Comme \(a \notin \Phi(x)\),
    la deuxième option est exclue, donc \(a \land x = 0\). Alors
    \(a = a \land 1 = a \land (x \lor \neg x) = (a \land x) \lor (a \land \neg x) = 0 \lor (a \land \neg x) = a \land \neg x\).
    D'où \(a \le \neg x\), ce qui signifie \(a \in \Phi(\neg x)\).
    L'égalité est établie.
  \end{itemize}
\end{enumerate}

\begin{center}\rule{0.5\linewidth}{0.5pt}\end{center}

\hypertarget{question-5-conclusion-et-cardinalituxe9}{%
\subsubsection{Question 5 : Conclusion et
cardinalité}\label{question-5-conclusion-et-cardinalituxe9}}

\begin{itemize}
\tightlist
\item
  L'application \(\Phi\) est injective (Q3) et préserve la structure
  d'algèbre de Boole (Q4).
\item
  De plus, pour tout sous-ensemble d'atomes
  \(S \subseteq \mathcal{A}t(B)\), l'élément
  \(x_S = \bigvee_{a \in S} a\) vérifie \(\Phi(x_S) = S\). L'application
  \(\Phi\) est donc surjective.
\item
  Par conséquent, \(\Phi\) est un \textbf{isomorphisme d'algèbres de
  Boole}.
\item
  Toute algèbre de Boole finie \(B\) est isomorphe à l'algèbre des
  parties \(\mathcal{P}(\mathcal{A}t(B))\).
\item
  Comme l'ensemble \(\mathcal{A}t(B)\) est fini (soit \(n\) son
  cardinal), le cardinal de \(\mathcal{P}(\mathcal{A}t(B))\) est égal à
  \(2^n\).
\end{itemize}

On en déduit le théorème fondamental : \textbf{Le cardinal de toute
algèbre de Boole finie est une puissance de \(2\)}.
\n
\hypertarget{exercice-10-thuxe9oruxe8me-de-compacituxe9-et-coloriage-de-graphes-infinis-niveau-ens}{%
\section{Exercice 10 : Théorème de compacité et coloriage de graphes
infinis (Niveau
ENS)}\label{exercice-10-thuxe9oruxe8me-de-compacituxe9-et-coloriage-de-graphes-infinis-niveau-ens}}

\hypertarget{uxe9noncuxe9}{%
\subsection{Énoncé}\label{uxe9noncuxe9}}

Soit \(G = (V, E)\) un graphe non orienté, où l'ensemble des sommets
\(V\) est infini (dénombrable ou non) et
\(E \subseteq \{\{u, v\} \mid u, v \in V, u \neq v\}\) est l'ensemble
des arêtes. Soit \(k \ge 1\) un entier naturel. On rappelle qu'un
\textbf{\(k\)-coloriage} propre de \(G\) est une application
\(f : V \to \{1, \dots, k\}\) telle que pour toute arête
\(\{u, v\} \in E\) : \[f(u) \neq f(v)\]

On suppose que le graphe \(G\) vérifie la propriété locale suivante :
\textgreater{} \textbf{Tout sous-graphe fini de \(G\) est
\(k\)-colorable.}

Le but de cet exercice est de démontrer, à l'aide du \textbf{théorème de
compacité de la logique propositionnelle}, le \textbf{théorème de De
Bruijn-Erdős} : \textgreater{} \textbf{Le graphe infini \(G\) est
lui-même \(k\)-colorable.}

\begin{enumerate}
\def\labelenumi{\arabic{enumi}.}
\tightlist
\item
  Modéliser le problème en définissant un ensemble de variables
  propositionnelles approprié.
\item
  Écrire les formules logiques modélisant les trois contraintes
  suivantes :

  \begin{itemize}
  \tightlist
  \item
    Contrainte 1 : Tout sommet possède au moins une couleur.
  \item
    Contrainte 2 : Tout sommet possède au plus une couleur.
  \item
    Contrainte 3 : Deux sommets adjacents n'ont pas la même couleur. On
    notera \(\Sigma\) l'ensemble de toutes ces clauses.
  \end{itemize}
\item
  Soit \(\Sigma_0\) un sous-ensemble fini de \(\Sigma\). Démontrer que
  \(\Sigma_0\) est satisfaisable.
\item
  Conclure par le théorème de compacité de la logique propositionnelle.
\end{enumerate}

\begin{center}\rule{0.5\linewidth}{0.5pt}\end{center}

\hypertarget{correction-duxe9tailluxe9e}{%
\subsection{Correction Détaillée}\label{correction-duxe9tailluxe9e}}

\hypertarget{question-1-choix-des-variables-propositionnelles}{%
\subsubsection{Question 1 : Choix des variables
propositionnelles}\label{question-1-choix-des-variables-propositionnelles}}

Pour chaque sommet \(i \in V\) et chaque couleur
\(c \in \{1, \dots, k\}\), nous définissons la variable propositionnelle
\(x_{i, c}\) dont la sémantique est :
\[x_{i, c} \equiv \text{« Le sommet } i \text{ est coloré avec la couleur } c \text{ »}\]
L'ensemble de nos variables propositionnelles est donc :
\[\mathcal{P} = \{x_{i, c} \mid i \in V, c \in \{1, \dots, k\}\}\]

\begin{center}\rule{0.5\linewidth}{0.5pt}\end{center}

\hypertarget{question-2-moduxe9lisation-des-contraintes-logiques}{%
\subsubsection{Question 2 : Modélisation des contraintes
logiques}\label{question-2-moduxe9lisation-des-contraintes-logiques}}

\begin{enumerate}
\def\labelenumi{\arabic{enumi}.}
\item
  \textbf{Chaque sommet a au moins une couleur :} Pour tout sommet
  \(i \in V\), la formule traduisant cette condition est une clause
  disjonctive sur toutes les couleurs disponibles :
  \[\Phi_i = \bigvee_{c=1}^k x_{i, c}\] Désignons par
  \(\Sigma_1 = \{\Phi_i \mid i \in V\}\) l'ensemble de ces clauses.
\item
  \textbf{Chaque sommet a au plus une couleur :} Pour tout sommet
  \(i \in V\) et pour chaque paire de couleurs distinctes
  \(c_1 \neq c_2\), le sommet ne peut pas porter les deux couleurs
  simultanément, ce qui s'écrit
  \(\neg(x_{i, c_1} \land x_{i, c_2}) \equiv \neg x_{i, c_1} \lor \neg x_{i, c_2}\).
  Désignons par \(\Sigma_2\) l'ensemble de toutes ces clauses
  d'exclusion :
  \[\Sigma_2 = \{\neg x_{i, c_1} \lor \neg x_{i, c_2} \mid i \in V, \quad 1 \le c_1 < c_2 \le k\}\]
\item
  \textbf{Les sommets adjacents ont des couleurs différentes :} Pour
  chaque arête \(\{i, j\} \in E\) et chaque couleur
  \(c \in \{1, \dots, k\}\), les sommets \(i\) et \(j\) ne peuvent pas
  être colorés tous deux en \(c\), ce qui s'écrit
  \(\neg(x_{i, c} \land x_{j, c}) \equiv \neg x_{i, c} \lor \neg x_{j, c}\).
  Désignons par \(\Sigma_3\) l'ensemble de ces clauses de conflit :
  \[\Sigma_3 = \{\neg x_{i, c} \lor \neg x_{j, c} \mid \{i, j\} \in E, \quad c \in \{1, \dots, k\}\}\]
\end{enumerate}

L'ensemble total de nos formules est
\(\Sigma = \Sigma_1 \cup \Sigma_2 \cup \Sigma_3\). Une valuation \(v\)
satisfaisant \(\Sigma\) correspond bijectivement à un \(k\)-coloriage
propre de \(G\).

\begin{center}\rule{0.5\linewidth}{0.5pt}\end{center}

\hypertarget{question-3-satisfaisabilituxe9-de-tout-sous-ensemble-fini-sigma_0}{%
\subsubsection{\texorpdfstring{Question 3 : Satisfaisabilité de tout
sous-ensemble fini
\(\Sigma_0\)}{Question 3 : Satisfaisabilité de tout sous-ensemble fini \textbackslash Sigma\_0}}\label{question-3-satisfaisabilituxe9-de-tout-sous-ensemble-fini-sigma_0}}

Soit \(\Sigma_0 \subseteq \Sigma\) un sous-ensemble fini de clauses. -
Puisque \(\Sigma_0\) est fini, il ne contient qu'un nombre fini de
clauses. - Chaque clause de \(\Sigma_0\) ne fait intervenir qu'un nombre
fini de sommets de \(V\). - Désignons par \(V_0 \subseteq V\) l'ensemble
des sommets apparaissant dans les clauses de \(\Sigma_0\). Puisque
\(\Sigma_0\) est fini et que chaque clause contient au plus \(k\)
sommets (les clauses de \(\Sigma_1\) ont \(1\) sommet, \(\Sigma_2\) ont
\(1\) sommet, et \(\Sigma_3\) ont \(2\) sommets), \(V_0\) est un
\textbf{ensemble fini}. - Soit \(G_0 = (V_0, E_0)\) le sous-graphe fini
de \(G\) induit par le sous-ensemble de sommets \(V_0\), où
\(E_0 = \{\{i, j\} \in E \mid i, j \in V_0\}\).

Par hypothèse de l'énoncé, le sous-graphe fini \(G_0\) est
\(k\)-colorable. Il existe donc une fonction de coloriage propre
\(f_0 : V_0 \to \{1, \dots, k\}\) sur ce sous-graphe. Définissons une
valuation \(v_0\) pour les variables propositionnelles apparaissant dans
\(\Sigma_0\) par :
\[v_0(x_{i, c}) = 1 \iff i \in V_0 \text{ et } f_0(i) = c\]

Vérifions que \(v_0\) satisfait toutes les clauses de \(\Sigma_0\) : -
Si une clause \(C \in \Sigma_0\) provient de \(\Sigma_1\), elle est de
la forme \(\bigvee_{c=1}^k x_{i, c}\) pour un certain \(i \in V_0\).
Comme \(f_0(i) \in \{1, \dots, k\}\), il existe un unique \(c^*\) tel
que \(f_0(i) = c^*\). Alors \(v_0(x_{i, c^*}) = 1\), la clause est donc
satisfaite. - Si \(C \in \Sigma_0\) provient de \(\Sigma_2\), elle est
de la forme \(\neg x_{i, c_1} \lor \neg x_{i, c_2}\) pour \(c_1 < c_2\).
Comme \(f_0(i)\) ne peut pas être égal à la fois à \(c_1\) et \(c_2\),
au moins l'une des variables \(x_{i, c_1}\) ou \(x_{i, c_2}\) est
évaluée à \(0\) par \(v_0\). La disjonction des négations est donc
évaluée à \(1\). La clause est satisfaite. - Si \(C \in \Sigma_0\)
provient de \(\Sigma_3\), elle est de la forme
\(\neg x_{i, c} \lor \neg x_{j, c}\) avec \(\{i, j\} \in E_0\). Puisque
\(f_0\) est un coloriage propre de \(G_0\), \(f_0(i) \neq f_0(j)\). Donc
\(i\) et \(j\) ne peuvent pas avoir tous deux la couleur \(c\). L'une
des variables \(x_{i, c}\) ou \(x_{j, c}\) vaut nécessairement \(0\), ce
qui satisfait la clause.

Puisque toutes les clauses de \(\Sigma_0\) sont satisfaites par \(v_0\),
le sous-ensemble fini \(\Sigma_0\) est \textbf{satisfaisable}.

\begin{center}\rule{0.5\linewidth}{0.5pt}\end{center}

\hypertarget{question-4-conclusion-par-compacituxe9}{%
\subsubsection{Question 4 : Conclusion par
compacité}\label{question-4-conclusion-par-compacituxe9}}

Nous venons de démontrer que pour tout sous-ensemble fini
\(\Sigma_0 \subseteq \Sigma\), \(\Sigma_0\) est satisfaisable. Par le
\textbf{théorème de compacité de la logique propositionnelle},
l'ensemble infini complet \(\Sigma\) est satisfaisable.

Il existe donc une valuation globale \(v^* : \mathcal{P} \to \{0, 1\}\)
telle que pour toute clause \(C \in \Sigma\), \(v^*(C) = 1\).
Définissons l'application de coloriage global
\(f^* : V \to \{1, \dots, k\}\) par :
\[f^*(i) = c \iff v^*(x_{i, c}) = 1\]

Vérifions que \(f^*\) est bien définie et est un coloriage propre : -
Les clauses de \(\Sigma_1\) garantissent que pour chaque \(i \in V\), il
existe au moins un \(c\) tel que \(v^*(x_{i, c}) = 1\). - Les clauses de
\(\Sigma_2\) garantissent que pour chaque \(i \in V\), il existe au plus
un \(c\) tel que \(v^*(x_{i, c}) = 1\). - Donc, \(f^*\) est une
application bien définie de \(V\) dans \(\{1, \dots, k\}\). - Les
clauses de \(\Sigma_3\) garantissent que si \(\{i, j\} \in E\), alors
pour tout \(c\), \(v^*(x_{i, c}) = 0\) ou \(v^*(x_{j, c}) = 0\). Donc
\(f^*(i) \neq f^*(j)\). Le coloriage est propre.

Le graphe infini \(G\) est donc \(k\)-colorable. \(\blacksquare\)
\n


\part{Travaux pratiques et simulations algorithmiques}
\hypertarget{tp-1-arbre-syntaxique-abstrait-pour-formules-logiques}{%
\section{TP 1 : Arbre syntaxique abstrait pour formules
logiques}\label{tp-1-arbre-syntaxique-abstrait-pour-formules-logiques}}

\hypertarget{objectif}{%
\subsection{Objectif}\label{objectif}}

L'objectif de ce TP est de concevoir et d'implémenter en Python une
structure de données d'\textbf{Arbre Syntaxique Abstrait (AST)}
permettant de modéliser, de représenter sous forme textuelle propre et
d'évaluer sémantiquement n'importe quelle formule bien formée de la
logique propositionnelle.

\begin{center}\rule{0.5\linewidth}{0.5pt}\end{center}

\hypertarget{conception-thuxe9orique}{%
\subsection{Conception Théorique}\label{conception-thuxe9orique}}

Pour modéliser une formule logique en programmation orientée objet, nous
utilisons le patron de conception \textbf{Composite}. - Une classe de
base abstraite \texttt{Formula} définit l'interface commune. - Les
variables propositionnelles sont représentées par des nœuds feuilles
(\texttt{Variable}). - Les connecteurs logiques correspondent à des
nœuds internes unaires (\texttt{Not}) ou binaires (\texttt{And},
\texttt{Or}, \texttt{Implies}, \texttt{Equiv}).

Chaque classe doit implémenter trois opérations clés : 1.
\texttt{\_\_str\_\_(self)} : retourne une représentation textuelle
parenthésée de la formule. 2. \texttt{get\_variables(self)} : retourne
l'ensemble (\texttt{set}) des variables propositionnelles présentes dans
la formule. 3. \texttt{evaluate(self,\ valuation)} : évalue la valeur de
vérité (\(True\) ou \(False\)) de la formule pour une valuation donnée
sous forme de dictionnaire Python \texttt{\{nom\_variable:\ booléen\}}.

\begin{center}\rule{0.5\linewidth}{0.5pt}\end{center}

\hypertarget{code-source-python-from-scratch}{%
\subsection{Code Source Python (from
scratch)}\label{code-source-python-from-scratch}}

Voici l'implémentation complète des classes en Python :

\begin{lstlisting}
\ImportTok{from}\NormalTok{ abc }\ImportTok{import}\NormalTok{ ABC, abstractmethod}

\KeywordTok{class}\NormalTok{ Formula(ABC):}
    \CommentTok{"""Classe abstraite de base pour toutes les formules propositionnelles."""}

    \AttributeTok{@abstractmethod}
    \KeywordTok{def} \FunctionTok{\_\_str\_\_}\NormalTok{(}\VariableTok{self}\NormalTok{):}
        \CommentTok{"""Retourne la représentation textuelle sous forme de chaîne de caractères."""}
        \ControlFlowTok{pass}

    \AttributeTok{@abstractmethod}
    \KeywordTok{def}\NormalTok{ get\_variables(}\VariableTok{self}\NormalTok{):}
        \CommentTok{"""Retourne un ensemble (set) de toutes les variables propositionnelles de la formule."""}
        \ControlFlowTok{pass}

    \AttributeTok{@abstractmethod}
    \KeywordTok{def}\NormalTok{ evaluate(}\VariableTok{self}\NormalTok{, valuation):}
        \CommentTok{"""Évalue la formule pour une valuation donnée (dictionnaire de type dict[str, bool])."""}
        \ControlFlowTok{pass}


\KeywordTok{class}\NormalTok{ Variable(Formula):}
    \CommentTok{"""Représente une variable propositionnelle (atome, feuille de l\textquotesingle{}arbre)."""}

    \KeywordTok{def} \FunctionTok{\_\_init\_\_}\NormalTok{(}\VariableTok{self}\NormalTok{, name):}
        \VariableTok{self}\NormalTok{.name }\OperatorTok{=}\NormalTok{ name}

    \KeywordTok{def} \FunctionTok{\_\_str\_\_}\NormalTok{(}\VariableTok{self}\NormalTok{):}
        \ControlFlowTok{return} \VariableTok{self}\NormalTok{.name}

    \KeywordTok{def}\NormalTok{ get\_variables(}\VariableTok{self}\NormalTok{):}
        \ControlFlowTok{return}\NormalTok{ \{}\VariableTok{self}\NormalTok{.name\}}

    \KeywordTok{def}\NormalTok{ evaluate(}\VariableTok{self}\NormalTok{, valuation):}
        \ControlFlowTok{if} \VariableTok{self}\NormalTok{.name }\KeywordTok{not} \KeywordTok{in}\NormalTok{ valuation:}
            \ControlFlowTok{raise} \PreprocessorTok{ValueError}\NormalTok{(}\SpecialStringTok{f"La variable }\SpecialCharTok{\{}\VariableTok{self}\SpecialCharTok{.}\NormalTok{name}\SpecialCharTok{\}}\SpecialStringTok{ n\textquotesingle{}est pas définie dans la valuation fournie."}\NormalTok{)}
        \ControlFlowTok{return}\NormalTok{ valuation[}\VariableTok{self}\NormalTok{.name]}


\KeywordTok{class}\NormalTok{ Not(Formula):}
    \CommentTok{"""Représente le connecteur de négation (noeud interne unaire)."""}

    \KeywordTok{def} \FunctionTok{\_\_init\_\_}\NormalTok{(}\VariableTok{self}\NormalTok{, child):}
        \VariableTok{self}\NormalTok{.child }\OperatorTok{=}\NormalTok{ child}

    \KeywordTok{def} \FunctionTok{\_\_str\_\_}\NormalTok{(}\VariableTok{self}\NormalTok{):}
        \ControlFlowTok{return} \SpecialStringTok{f"\textasciitilde{}(}\SpecialCharTok{\{}\VariableTok{self}\SpecialCharTok{.}\NormalTok{child}\SpecialCharTok{\}}\SpecialStringTok{)"}

    \KeywordTok{def}\NormalTok{ get\_variables(}\VariableTok{self}\NormalTok{):}
        \ControlFlowTok{return} \VariableTok{self}\NormalTok{.child.get\_variables()}

    \KeywordTok{def}\NormalTok{ evaluate(}\VariableTok{self}\NormalTok{, valuation):}
        \ControlFlowTok{return} \KeywordTok{not} \VariableTok{self}\NormalTok{.child.evaluate(valuation)}


\KeywordTok{class}\NormalTok{ BinaryOperator(Formula, ABC):}
    \CommentTok{"""Classe intermédiaire pour les opérateurs binaires (conjonction, disjonction, etc.)."""}

    \KeywordTok{def} \FunctionTok{\_\_init\_\_}\NormalTok{(}\VariableTok{self}\NormalTok{, left, right):}
        \VariableTok{self}\NormalTok{.left }\OperatorTok{=}\NormalTok{ left}
        \VariableTok{self}\NormalTok{.right }\OperatorTok{=}\NormalTok{ right}

    \KeywordTok{def}\NormalTok{ get\_variables(}\VariableTok{self}\NormalTok{):}
        \ControlFlowTok{return} \VariableTok{self}\NormalTok{.left.get\_variables().union(}\VariableTok{self}\NormalTok{.right.get\_variables())}


\KeywordTok{class}\NormalTok{ And(BinaryOperator):}
    \CommentTok{"""Représente le connecteur de conjonction ET (noeud interne binaire)."""}

    \KeywordTok{def} \FunctionTok{\_\_str\_\_}\NormalTok{(}\VariableTok{self}\NormalTok{):}
        \ControlFlowTok{return} \SpecialStringTok{f"(}\SpecialCharTok{\{}\VariableTok{self}\SpecialCharTok{.}\NormalTok{left}\SpecialCharTok{\}}\SpecialStringTok{ \& }\SpecialCharTok{\{}\VariableTok{self}\SpecialCharTok{.}\NormalTok{right}\SpecialCharTok{\}}\SpecialStringTok{)"}

    \KeywordTok{def}\NormalTok{ evaluate(}\VariableTok{self}\NormalTok{, valuation):}
        \ControlFlowTok{return} \VariableTok{self}\NormalTok{.left.evaluate(valuation) }\KeywordTok{and} \VariableTok{self}\NormalTok{.right.evaluate(valuation)}


\KeywordTok{class}\NormalTok{ Or(BinaryOperator):}
    \CommentTok{"""Représente le connecteur de disjonction OU (noeud interne binaire)."""}

    \KeywordTok{def} \FunctionTok{\_\_str\_\_}\NormalTok{(}\VariableTok{self}\NormalTok{):}
        \ControlFlowTok{return} \SpecialStringTok{f"(}\SpecialCharTok{\{}\VariableTok{self}\SpecialCharTok{.}\NormalTok{left}\SpecialCharTok{\}}\SpecialStringTok{ | }\SpecialCharTok{\{}\VariableTok{self}\SpecialCharTok{.}\NormalTok{right}\SpecialCharTok{\}}\SpecialStringTok{)"}

    \KeywordTok{def}\NormalTok{ evaluate(}\VariableTok{self}\NormalTok{, valuation):}
        \ControlFlowTok{return} \VariableTok{self}\NormalTok{.left.evaluate(valuation) }\KeywordTok{or} \VariableTok{self}\NormalTok{.right.evaluate(valuation)}


\KeywordTok{class}\NormalTok{ Implies(BinaryOperator):}
    \CommentTok{"""Représente le connecteur d\textquotesingle{}implication (SI ... ALORS ...)."""}

    \KeywordTok{def} \FunctionTok{\_\_str\_\_}\NormalTok{(}\VariableTok{self}\NormalTok{):}
        \ControlFlowTok{return} \SpecialStringTok{f"(}\SpecialCharTok{\{}\VariableTok{self}\SpecialCharTok{.}\NormalTok{left}\SpecialCharTok{\}}\SpecialStringTok{ =\textgreater{} }\SpecialCharTok{\{}\VariableTok{self}\SpecialCharTok{.}\NormalTok{right}\SpecialCharTok{\}}\SpecialStringTok{)"}

    \KeywordTok{def}\NormalTok{ evaluate(}\VariableTok{self}\NormalTok{, valuation):}
        \CommentTok{\# A =\textgreater{} B est équivalent à (non A) ou B}
        \ControlFlowTok{return}\NormalTok{ (}\KeywordTok{not} \VariableTok{self}\NormalTok{.left.evaluate(valuation)) }\KeywordTok{or} \VariableTok{self}\NormalTok{.right.evaluate(valuation)}


\KeywordTok{class}\NormalTok{ Equiv(BinaryOperator):}
    \CommentTok{"""Représente le connecteur d\textquotesingle{}équivalence équivalent (SI ET SEULEMENT SI)."""}

    \KeywordTok{def} \FunctionTok{\_\_str\_\_}\NormalTok{(}\VariableTok{self}\NormalTok{):}
        \ControlFlowTok{return} \SpecialStringTok{f"(}\SpecialCharTok{\{}\VariableTok{self}\SpecialCharTok{.}\NormalTok{left}\SpecialCharTok{\}}\SpecialStringTok{ \textless{}=\textgreater{} }\SpecialCharTok{\{}\VariableTok{self}\SpecialCharTok{.}\NormalTok{right}\SpecialCharTok{\}}\SpecialStringTok{)"}

    \KeywordTok{def}\NormalTok{ evaluate(}\VariableTok{self}\NormalTok{, valuation):}
        \ControlFlowTok{return} \VariableTok{self}\NormalTok{.left.evaluate(valuation) }\OperatorTok{==} \VariableTok{self}\NormalTok{.right.evaluate(valuation)}
\end{lstlisting}

\begin{center}\rule{0.5\linewidth}{0.5pt}\end{center}

\hypertarget{validation-et-exemple-dutilisation}{%
\subsection{Validation et Exemple
d'utilisation}\label{validation-et-exemple-dutilisation}}

Ajoutons un script de test pour valider notre implémentation :

\begin{lstlisting}
\ControlFlowTok{if} \VariableTok{\_\_name\_\_} \OperatorTok{==} \StringTok{"\_\_main\_\_"}\NormalTok{:}
    \CommentTok{\# Définition des variables}
\NormalTok{    p }\OperatorTok{=}\NormalTok{ Variable(}\StringTok{"P"}\NormalTok{)}
\NormalTok{    q }\OperatorTok{=}\NormalTok{ Variable(}\StringTok{"Q"}\NormalTok{)}
\NormalTok{    r }\OperatorTok{=}\NormalTok{ Variable(}\StringTok{"R"}\NormalTok{)}

    \CommentTok{\# Construction de la formule : (P \& Q) =\textgreater{} (R | \textasciitilde{}P)}
\NormalTok{    formule }\OperatorTok{=}\NormalTok{ Implies(And(p, q), Or(r, Not(p)))}

    \BuiltInTok{print}\NormalTok{(}\StringTok{"Formule logique générée :"}\NormalTok{)}
    \BuiltInTok{print}\NormalTok{(formule)  }\CommentTok{\# Doit afficher : ((P \& Q) =\textgreater{} (R | \textasciitilde{}(P)))}

    \BuiltInTok{print}\NormalTok{(}\StringTok{"}\CharTok{\textbackslash{}n}\StringTok{Extraction des variables :"}\NormalTok{)}
    \BuiltInTok{print}\NormalTok{(formule.get\_variables())  }\CommentTok{\# Doit afficher : \{\textquotesingle{}P\textquotesingle{}, \textquotesingle{}Q\textquotesingle{}, \textquotesingle{}R\textquotesingle{}\}}

    \CommentTok{\# Définition d\textquotesingle{}une valuation}
\NormalTok{    valuation\_1 }\OperatorTok{=}\NormalTok{ \{}\StringTok{"P"}\NormalTok{: }\VariableTok{True}\NormalTok{, }\StringTok{"Q"}\NormalTok{: }\VariableTok{True}\NormalTok{, }\StringTok{"R"}\NormalTok{: }\VariableTok{False}\NormalTok{\}}
    \BuiltInTok{print}\NormalTok{(}\SpecialStringTok{f"}\CharTok{\textbackslash{}n}\SpecialStringTok{Évaluation avec la valuation }\SpecialCharTok{\{}\NormalTok{valuation\_1}\SpecialCharTok{\}}\SpecialStringTok{ :"}\NormalTok{)}
\NormalTok{    res\_1 }\OperatorTok{=}\NormalTok{ formule.evaluate(valuation\_1)}
    \BuiltInTok{print}\NormalTok{(}\SpecialStringTok{f"Résultat : }\SpecialCharTok{\{}\NormalTok{res\_1}\SpecialCharTok{\}}\SpecialStringTok{"}\NormalTok{)  }\CommentTok{\# Doit afficher : False (car Vrai =\textgreater{} Faux est Faux)}

\NormalTok{    valuation\_2 }\OperatorTok{=}\NormalTok{ \{}\StringTok{"P"}\NormalTok{: }\VariableTok{False}\NormalTok{, }\StringTok{"Q"}\NormalTok{: }\VariableTok{True}\NormalTok{, }\StringTok{"R"}\NormalTok{: }\VariableTok{False}\NormalTok{\}}
    \BuiltInTok{print}\NormalTok{(}\SpecialStringTok{f"}\CharTok{\textbackslash{}n}\SpecialStringTok{Évaluation avec la valuation }\SpecialCharTok{\{}\NormalTok{valuation\_2}\SpecialCharTok{\}}\SpecialStringTok{ :"}\NormalTok{)}
\NormalTok{    res\_2 }\OperatorTok{=}\NormalTok{ formule.evaluate(valuation\_2)}
    \BuiltInTok{print}\NormalTok{(}\SpecialStringTok{f"Résultat : }\SpecialCharTok{\{}\NormalTok{res\_2}\SpecialCharTok{\}}\SpecialStringTok{"}\NormalTok{)  }\CommentTok{\# Doit afficher : True (car Faux =\textgreater{} Vrai est Vrai)}
\end{lstlisting}
\n
\hypertarget{tp-2-guxe9nuxe9rateur-automatique-de-tables-de-vuxe9rituxe9}{%
\section{TP 2 : Générateur automatique de tables de
vérité}\label{tp-2-guxe9nuxe9rateur-automatique-de-tables-de-vuxe9rituxe9}}

\hypertarget{objectif}{%
\subsection{Objectif}\label{objectif}}

L'objectif de ce TP est de concevoir un générateur automatique de tables
de vérité en Python. À partir de n'importe quel arbre de formule logique
implémenté au TP 1, le programme doit lister toutes les variables
libres, générer les \(2^n\) combinaisons de valeurs de vérité possibles,
évaluer la formule et imprimer la table de vérité finale proprement
formatée en Markdown.

\begin{center}\rule{0.5\linewidth}{0.5pt}\end{center}

\hypertarget{algorithme-et-muxe9thode}{%
\subsection{Algorithme et Méthode}\label{algorithme-et-muxe9thode}}

\begin{enumerate}
\def\labelenumi{\arabic{enumi}.}
\tightlist
\item
  \textbf{Extraction et tri des variables :} On extrait l'ensemble des
  variables propositionnelles libres de la formule en appelant
  \texttt{formula.get\_variables()}. On convertit cet ensemble en une
  liste ordonnée alphabétiquement afin de garantir un affichage
  reproductible et cohérent des colonnes.
\item
  \textbf{Génération de l'espace des valuations \(\{0, 1\}^n\) :} On
  utilise la récursion pour engendrer toutes les \(2^n\) configurations
  de vérité sans utiliser de bibliothèque standard afin de rester
  strictement ``from scratch''.
\item
  \textbf{Évaluation et formatage :} Pour chaque configuration
  (valuation), on évalue la formule à l'aide de sa méthode
  \texttt{.evaluate(valuation)}. On affiche chaque ligne avec des
  symboles lisibles (par exemple \texttt{1} ou \texttt{T} pour True,
  \texttt{0} ou \texttt{F} pour False) et on trace la grille de la table
  au format Markdown.
\end{enumerate}

\begin{center}\rule{0.5\linewidth}{0.5pt}\end{center}

\hypertarget{code-source-python}{%
\subsection{Code Source Python}\label{code-source-python}}

\begin{lstlisting}
\CommentTok{\# Import des classes du TP 1}
\CommentTok{\# On suppose les classes du TP 1 présentes dans le même fichier ou importées}
\ImportTok{from}\NormalTok{ tp\_01\_ast }\ImportTok{import}\NormalTok{ Formula, Variable, Not, And, Or, Implies, Equiv}

\KeywordTok{def}\NormalTok{ generate\_combinations(n):}
    \CommentTok{"""Génère récursivement toutes les 2\^{}n combinaisons de valeurs de vérité de True à False."""}
    \ControlFlowTok{if}\NormalTok{ n }\OperatorTok{==} \DecValTok{0}\NormalTok{:}
        \ControlFlowTok{return}\NormalTok{ [()]}
\NormalTok{    sub\_combinations }\OperatorTok{=}\NormalTok{ generate\_combinations(n }\OperatorTok{{-}} \DecValTok{1}\NormalTok{)}
    \ControlFlowTok{return}\NormalTok{ [(}\VariableTok{True}\NormalTok{,) }\OperatorTok{+}\NormalTok{ combo }\ControlFlowTok{for}\NormalTok{ combo }\KeywordTok{in}\NormalTok{ sub\_combinations] }\OperatorTok{+}\NormalTok{ [(}\VariableTok{False}\NormalTok{,) }\OperatorTok{+}\NormalTok{ combo }\ControlFlowTok{for}\NormalTok{ combo }\KeywordTok{in}\NormalTok{ sub\_combinations]}

\KeywordTok{def}\NormalTok{ generate\_truth\_table(formula: Formula):}
    \CommentTok{"""Génère et affiche la table de vérité d\textquotesingle{}une formule au format Markdown."""}

    \CommentTok{\# 1. Extraction et tri des variables propositionnelles}
\NormalTok{    vars\_list }\OperatorTok{=} \BuiltInTok{sorted}\NormalTok{(}\BuiltInTok{list}\NormalTok{(formula.get\_variables()))}
\NormalTok{    n }\OperatorTok{=} \BuiltInTok{len}\NormalTok{(vars\_list)}

    \CommentTok{\# 2. Construction des en{-}têtes Markdown}
\NormalTok{    headers }\OperatorTok{=}\NormalTok{ vars\_list }\OperatorTok{+}\NormalTok{ [}\BuiltInTok{str}\NormalTok{(formula)]}
\NormalTok{    header\_line }\OperatorTok{=} \StringTok{" | "}\NormalTok{.join(headers)}
\NormalTok{    separator\_line }\OperatorTok{=} \StringTok{" | "}\NormalTok{.join([}\StringTok{"{-}{-}{-}"}\NormalTok{] }\OperatorTok{*} \BuiltInTok{len}\NormalTok{(headers))}

\NormalTok{    table\_lines }\OperatorTok{=}\NormalTok{ [header\_line, separator\_line]}

    \CommentTok{\# 3. Génération de toutes les valuations possibles (2\^{}n)}
\NormalTok{    combinations }\OperatorTok{=}\NormalTok{ generate\_combinations(n)}

    \CommentTok{\# Pour respecter l\textquotesingle{}ordre traditionnel (Vrai en premier, c\textquotesingle{}est{-}à{-}dire 1 puis 0),}
    \CommentTok{\# product avec [True, False] trie naturellement de True à False.}

\NormalTok{    tautology }\OperatorTok{=} \VariableTok{True}
\NormalTok{    contradiction }\OperatorTok{=} \VariableTok{True}

    \ControlFlowTok{for}\NormalTok{ combo }\KeywordTok{in}\NormalTok{ combinations:}
        \CommentTok{\# Création du dictionnaire de valuation}
\NormalTok{        valuation }\OperatorTok{=}\NormalTok{ \{vars\_list[i]: combo[i] }\ControlFlowTok{for}\NormalTok{ i }\KeywordTok{in} \BuiltInTok{range}\NormalTok{(n)\}}

        \CommentTok{\# Évaluation de la formule pour cette valuation}
\NormalTok{        result }\OperatorTok{=}\NormalTok{ formula.evaluate(valuation)}

        \CommentTok{\# Suivi de la nature de la formule}
        \ControlFlowTok{if}\NormalTok{ result:}
\NormalTok{            contradiction }\OperatorTok{=} \VariableTok{False}
        \ControlFlowTok{else}\NormalTok{:}
\NormalTok{            tautology }\OperatorTok{=} \VariableTok{False}

        \CommentTok{\# Conversion des valeurs booléennes en chaînes de caractères (\textquotesingle{}1\textquotesingle{} ou \textquotesingle{}0\textquotesingle{})}
\NormalTok{        row\_values }\OperatorTok{=}\NormalTok{ [(}\StringTok{"1"} \ControlFlowTok{if}\NormalTok{ val }\ControlFlowTok{else} \StringTok{"0"}\NormalTok{) }\ControlFlowTok{for}\NormalTok{ val }\KeywordTok{in}\NormalTok{ combo] }\OperatorTok{+}\NormalTok{ [(}\StringTok{"1"} \ControlFlowTok{if}\NormalTok{ result }\ControlFlowTok{else} \StringTok{"0"}\NormalTok{)]}
\NormalTok{        table\_lines.append(}\StringTok{" | "}\NormalTok{.join(row\_values))}

    \CommentTok{\# Affichage de la table de vérité}
    \BuiltInTok{print}\NormalTok{(}\StringTok{"}\CharTok{\textbackslash{}n}\StringTok{"}\NormalTok{.join(table\_lines))}

    \CommentTok{\# 4. Diagnostic de la formule}
    \BuiltInTok{print}\NormalTok{(}\StringTok{"}\CharTok{\textbackslash{}n}\StringTok{{-}{-}{-}"}\NormalTok{)}
    \BuiltInTok{print}\NormalTok{(}\StringTok{"\#\#\# Diagnostic de la formule :"}\NormalTok{)}
    \ControlFlowTok{if}\NormalTok{ tautology:}
        \BuiltInTok{print}\NormalTok{(}\StringTok{"La formule est une **Tautologie** (toujours vraie)."}\NormalTok{)}
    \ControlFlowTok{elif}\NormalTok{ contradiction:}
        \BuiltInTok{print}\NormalTok{(}\StringTok{"La formule est une **Contradiction** (toujours fausse)."}\NormalTok{)}
    \ControlFlowTok{else}\NormalTok{:}
        \BuiltInTok{print}\NormalTok{(}\StringTok{"La formule est **Contingente** (satisfaisable, possède au moins un modèle et un contre{-}modèle)."}\NormalTok{)}
\end{lstlisting}

\begin{center}\rule{0.5\linewidth}{0.5pt}\end{center}

\hypertarget{exemple-dexuxe9cution}{%
\subsection{Exemple d'Exécution}\label{exemple-dexuxe9cution}}

Modélisons la formule du tiers exclu et l'implication logique :

\begin{lstlisting}
\ControlFlowTok{if} \VariableTok{\_\_name\_\_} \OperatorTok{==} \StringTok{"\_\_main\_\_"}\NormalTok{:}
    \CommentTok{\# Test 1 : Tiers exclu : P | \textasciitilde{}P}
\NormalTok{    p }\OperatorTok{=}\NormalTok{ Variable(}\StringTok{"P"}\NormalTok{)}
\NormalTok{    tiers\_exclu }\OperatorTok{=}\NormalTok{ Or(p, Not(p))}
    \BuiltInTok{print}\NormalTok{(}\StringTok{"\#\#\# Table de vérité de la formule : P | \textasciitilde{}P"}\NormalTok{)}
\NormalTok{    generate\_truth\_table(tiers\_exclu)}

    \CommentTok{\# Test 2 : Formule contingente : (P \& Q) =\textgreater{} R}
\NormalTok{    q }\OperatorTok{=}\NormalTok{ Variable(}\StringTok{"Q"}\NormalTok{)}
\NormalTok{    r }\OperatorTok{=}\NormalTok{ Variable(}\StringTok{"R"}\NormalTok{)}
\NormalTok{    contingence }\OperatorTok{=}\NormalTok{ Implies(And(p, q), r)}
    \BuiltInTok{print}\NormalTok{(}\StringTok{"}\CharTok{\textbackslash{}n}\StringTok{\#\#\# Table de vérité de la formule : (P \& Q) =\textgreater{} R"}\NormalTok{)}
\NormalTok{    generate\_truth\_table(contingence)}

    \CommentTok{\# Test 3 : Loi de Peirce : ((P =\textgreater{} Q) =\textgreater{} P) =\textgreater{} P}
\NormalTok{    peirce }\OperatorTok{=}\NormalTok{ Implies(Implies(Implies(p, q), p), p)}
    \BuiltInTok{print}\NormalTok{(}\StringTok{"}\CharTok{\textbackslash{}n}\StringTok{\#\#\# Table de vérité de la Loi de Peirce : ((P =\textgreater{} Q) =\textgreater{} P) =\textgreater{} P"}\NormalTok{)}
\NormalTok{    generate\_truth\_table(peirce)}
\end{lstlisting}

\hypertarget{rendu-attendu-en-console-pour-le-test-3-loi-de-peirce}{%
\subsubsection{Rendu attendu en console pour le Test 3 (Loi de Peirce)
:}\label{rendu-attendu-en-console-pour-le-test-3-loi-de-peirce}}

\begin{lstlisting}
\NormalTok{P | Q | (((P =\textgreater{} Q) =\textgreater{} P) =\textgreater{} P)}
\NormalTok{{-}{-}{-} | {-}{-}{-} | {-}{-}{-}}
\NormalTok{1 | 1 | 1}
\NormalTok{1 | 0 | 1}
\NormalTok{0 | 1 | 1}
\NormalTok{0 | 0 | 1}

\NormalTok{{-}{-}{-}}
\FunctionTok{\#\#\# Diagnostic de la formule :}
\NormalTok{La formule est une **Tautologie** (toujours vraie).}
\end{lstlisting}
\n
\hypertarget{tp-3-transformation-de-tseitin-pour-la-mise-en-cnf-linuxe9aire}{%
\section{TP 3 : Transformation de Tseitin pour la mise en CNF
linéaire}\label{tp-3-transformation-de-tseitin-pour-la-mise-en-cnf-linuxe9aire}}

\hypertarget{objectif}{%
\subsection{Objectif}\label{objectif}}

L'objectif de ce TP est d'implémenter la \textbf{transformation de
Tseitin} en Python. La conversion naïve d'une formule en Forme Normale
Conjonctive (CNF) par simple application des règles de distributivité
peut conduire à une explosion exponentielle de la taille de la formule.
La transformation de Tseitin contourne ce problème en introduisant des
variables auxiliaires pour chaque sous-formule, permettant de générer
une formule en CNF équisatisfiable de taille \textbf{linéaire} par
rapport à la formule initiale.

\begin{center}\rule{0.5\linewidth}{0.5pt}\end{center}

\hypertarget{principes-thuxe9oriques-et-moduxe9lisation}{%
\subsection{Principes Théoriques et
Modélisation}\label{principes-thuxe9oriques-et-moduxe9lisation}}

Pour chaque sous-formule (nœud interne ou feuille de notre AST), la
transformation de Tseitin introduit une nouvelle variable
propositionnelle \(x_i\). On ajoute ensuite des clauses qui forcent la
variable \(x_i\) à être équivalente à l'application de l'opérateur
logique sur ses enfants.

\hypertarget{uxe9quivalences-logiques-en-cnf}{%
\subsubsection{Équivalences logiques en CNF
:}\label{uxe9quivalences-logiques-en-cnf}}

\begin{enumerate}
\def\labelenumi{\arabic{enumi}.}
\tightlist
\item
  \textbf{Négation :} \(x_C \Leftrightarrow \neg x_A\)

  \begin{itemize}
  \tightlist
  \item
    Équivalent à : \((\neg x_C \lor \neg x_A) \land (x_C \lor x_A)\)
  \end{itemize}
\item
  \textbf{Conjonction :} \(x_C \Leftrightarrow (x_A \land x_B)\)

  \begin{itemize}
  \tightlist
  \item
    Équivalent à :
    \((\neg x_C \lor x_A) \land (\neg x_C \lor x_B) \land (\neg x_A \lor \neg x_B \lor x_C)\)
  \end{itemize}
\item
  \textbf{Disjonction :} \(x_C \Leftrightarrow (x_A \lor x_B)\)

  \begin{itemize}
  \tightlist
  \item
    Équivalent à :
    \((\neg x_C \lor x_A \lor x_B) \land (\neg x_A \lor x_C) \land (\neg x_B \lor x_C)\)
  \end{itemize}
\item
  \textbf{Implication :} \(x_C \Leftrightarrow (x_A \Rightarrow x_B)\)

  \begin{itemize}
  \tightlist
  \item
    Équivalent à :
    \((\neg x_C \lor \neg x_A \lor x_B) \land (x_A \lor x_C) \land (\neg x_B \lor x_C)\)
  \end{itemize}
\end{enumerate}

La CNF finale est constituée de la conjonction de toutes ces contraintes
de nœuds, à laquelle on ajoute la clause unitaire \((x_{root})\) qui
force la formule générale à être vraie.

\begin{center}\rule{0.5\linewidth}{0.5pt}\end{center}

\hypertarget{code-source-python}{%
\subsection{Code Source Python}\label{code-source-python}}

Chaque clause sera représentée sous forme d'un ensemble de littéraux (un
littéral étant une chaîne de caractères comme \texttt{"P"} ou
\texttt{"\textasciitilde{}P"}). La CNF est une liste de clauses (liste
de sets).

\begin{lstlisting}
\CommentTok{\# Import des classes du TP 1}
\ImportTok{from}\NormalTok{ tp\_01\_ast }\ImportTok{import}\NormalTok{ Formula, Variable, Not, And, Or, Implies, Equiv}

\KeywordTok{class}\NormalTok{ TseitinTransformer:}
    \KeywordTok{def} \FunctionTok{\_\_init\_\_}\NormalTok{(}\VariableTok{self}\NormalTok{):}
        \VariableTok{self}\NormalTok{.counter }\OperatorTok{=} \DecValTok{0}
        \VariableTok{self}\NormalTok{.clauses }\OperatorTok{=}\NormalTok{ []}

    \KeywordTok{def}\NormalTok{ \_new\_variable(}\VariableTok{self}\NormalTok{):}
        \CommentTok{"""Crée une variable auxiliaire unique (ex: \_t1, \_t2, ...)."""}
        \VariableTok{self}\NormalTok{.counter }\OperatorTok{+=} \DecValTok{1}
        \ControlFlowTok{return} \SpecialStringTok{f"\_t}\SpecialCharTok{\{}\VariableTok{self}\SpecialCharTok{.}\NormalTok{counter}\SpecialCharTok{\}}\SpecialStringTok{"}

    \KeywordTok{def}\NormalTok{ transform(}\VariableTok{self}\NormalTok{, formula: Formula):}
        \CommentTok{"""}
\CommentTok{        Point d\textquotesingle{}entrée principal. }
\CommentTok{        Retourne : (variable\_racine, liste\_de\_clauses\_CNF)}
\CommentTok{        """}
        \VariableTok{self}\NormalTok{.clauses }\OperatorTok{=}\NormalTok{ []}
        \VariableTok{self}\NormalTok{.counter }\OperatorTok{=} \DecValTok{0}
\NormalTok{        root\_var }\OperatorTok{=} \VariableTok{self}\NormalTok{.\_traverse(formula)}
        \CommentTok{\# On force la racine à être Vraie en ajoutant la clause unitaire \{root\_var\}}
        \VariableTok{self}\NormalTok{.clauses.append(\{root\_var\})}
        \ControlFlowTok{return}\NormalTok{ root\_var, }\VariableTok{self}\NormalTok{.clauses}

    \KeywordTok{def}\NormalTok{ \_traverse(}\VariableTok{self}\NormalTok{, node: Formula):}
        \CommentTok{"""Parcours récursif de l\textquotesingle{}AST."""}
        \ControlFlowTok{if} \BuiltInTok{isinstance}\NormalTok{(node, Variable):}
            \ControlFlowTok{return}\NormalTok{ node.name}

        \ControlFlowTok{elif} \BuiltInTok{isinstance}\NormalTok{(node, Not):}
\NormalTok{            child\_var }\OperatorTok{=} \VariableTok{self}\NormalTok{.\_traverse(node.child)}
\NormalTok{            fresh\_var }\OperatorTok{=} \VariableTok{self}\NormalTok{.\_new\_variable()}
            \CommentTok{\# Contrainte : fresh\_var \textless{}=\textgreater{} \textasciitilde{}child\_var}
            \CommentTok{\# (\textasciitilde{}fresh\_var v \textasciitilde{}child\_var) \& (fresh\_var v child\_var)}
            \VariableTok{self}\NormalTok{.clauses.append(\{}\SpecialStringTok{f"\textasciitilde{}}\SpecialCharTok{\{}\NormalTok{fresh\_var}\SpecialCharTok{\}}\SpecialStringTok{"}\NormalTok{, }\SpecialStringTok{f"\textasciitilde{}}\SpecialCharTok{\{}\NormalTok{child\_var}\SpecialCharTok{\}}\SpecialStringTok{"}\NormalTok{\})}
            \VariableTok{self}\NormalTok{.clauses.append(\{fresh\_var, child\_var\})}
            \ControlFlowTok{return}\NormalTok{ fresh\_var}

        \ControlFlowTok{elif} \BuiltInTok{isinstance}\NormalTok{(node, And):}
\NormalTok{            left\_var }\OperatorTok{=} \VariableTok{self}\NormalTok{.\_traverse(node.left)}
\NormalTok{            right\_var }\OperatorTok{=} \VariableTok{self}\NormalTok{.\_traverse(node.right)}
\NormalTok{            fresh\_var }\OperatorTok{=} \VariableTok{self}\NormalTok{.\_new\_variable()}
            \CommentTok{\# Contrainte : fresh\_var \textless{}=\textgreater{} (left\_var \& right\_var)}
            \CommentTok{\# (\textasciitilde{}fresh\_var v left\_var) \& (\textasciitilde{}fresh\_var v right\_var) \& (\textasciitilde{}left\_var v \textasciitilde{}right\_var v fresh\_var)}
            \VariableTok{self}\NormalTok{.clauses.append(\{}\SpecialStringTok{f"\textasciitilde{}}\SpecialCharTok{\{}\NormalTok{fresh\_var}\SpecialCharTok{\}}\SpecialStringTok{"}\NormalTok{, left\_var\})}
            \VariableTok{self}\NormalTok{.clauses.append(\{}\SpecialStringTok{f"\textasciitilde{}}\SpecialCharTok{\{}\NormalTok{fresh\_var}\SpecialCharTok{\}}\SpecialStringTok{"}\NormalTok{, right\_var\})}
            \VariableTok{self}\NormalTok{.clauses.append(\{}\SpecialStringTok{f"\textasciitilde{}}\SpecialCharTok{\{}\NormalTok{left\_var}\SpecialCharTok{\}}\SpecialStringTok{"}\NormalTok{, }\SpecialStringTok{f"\textasciitilde{}}\SpecialCharTok{\{}\NormalTok{right\_var}\SpecialCharTok{\}}\SpecialStringTok{"}\NormalTok{, fresh\_var\})}
            \ControlFlowTok{return}\NormalTok{ fresh\_var}

        \ControlFlowTok{elif} \BuiltInTok{isinstance}\NormalTok{(node, Or):}
\NormalTok{            left\_var }\OperatorTok{=} \VariableTok{self}\NormalTok{.\_traverse(node.left)}
\NormalTok{            right\_var }\OperatorTok{=} \VariableTok{self}\NormalTok{.\_traverse(node.right)}
\NormalTok{            fresh\_var }\OperatorTok{=} \VariableTok{self}\NormalTok{.\_new\_variable()}
            \CommentTok{\# Contrainte : fresh\_var \textless{}=\textgreater{} (left\_var | right\_var)}
            \CommentTok{\# (\textasciitilde{}fresh\_var v left\_var v right\_var) \& (\textasciitilde{}left\_var v fresh\_var) \& (\textasciitilde{}right\_var v fresh\_var)}
            \VariableTok{self}\NormalTok{.clauses.append(\{}\SpecialStringTok{f"\textasciitilde{}}\SpecialCharTok{\{}\NormalTok{fresh\_var}\SpecialCharTok{\}}\SpecialStringTok{"}\NormalTok{, left\_var, right\_var\})}
            \VariableTok{self}\NormalTok{.clauses.append(\{}\SpecialStringTok{f"\textasciitilde{}}\SpecialCharTok{\{}\NormalTok{left\_var}\SpecialCharTok{\}}\SpecialStringTok{"}\NormalTok{, fresh\_var\})}
            \VariableTok{self}\NormalTok{.clauses.append(\{}\SpecialStringTok{f"\textasciitilde{}}\SpecialCharTok{\{}\NormalTok{right\_var}\SpecialCharTok{\}}\SpecialStringTok{"}\NormalTok{, fresh\_var\})}
            \ControlFlowTok{return}\NormalTok{ fresh\_var}

        \ControlFlowTok{elif} \BuiltInTok{isinstance}\NormalTok{(node, Implies):}
\NormalTok{            left\_var }\OperatorTok{=} \VariableTok{self}\NormalTok{.\_traverse(node.left)}
\NormalTok{            right\_var }\OperatorTok{=} \VariableTok{self}\NormalTok{.\_traverse(node.right)}
\NormalTok{            fresh\_var }\OperatorTok{=} \VariableTok{self}\NormalTok{.\_new\_variable()}
            \CommentTok{\# Contrainte : fresh\_var \textless{}=\textgreater{} (left\_var =\textgreater{} right\_var)}
            \CommentTok{\# equivalent a : fresh\_var \textless{}=\textgreater{} (\textasciitilde{}left\_var | right\_var)}
            \CommentTok{\# (\textasciitilde{}fresh\_var v \textasciitilde{}left\_var v right\_var) \& (left\_var v fresh\_var) \& (\textasciitilde{}right\_var v fresh\_var)}
            \VariableTok{self}\NormalTok{.clauses.append(\{}\SpecialStringTok{f"\textasciitilde{}}\SpecialCharTok{\{}\NormalTok{fresh\_var}\SpecialCharTok{\}}\SpecialStringTok{"}\NormalTok{, }\SpecialStringTok{f"\textasciitilde{}}\SpecialCharTok{\{}\NormalTok{left\_var}\SpecialCharTok{\}}\SpecialStringTok{"}\NormalTok{, right\_var\})}
            \VariableTok{self}\NormalTok{.clauses.append(\{left\_var, fresh\_var\})}
            \VariableTok{self}\NormalTok{.clauses.append(\{}\SpecialStringTok{f"\textasciitilde{}}\SpecialCharTok{\{}\NormalTok{right\_var}\SpecialCharTok{\}}\SpecialStringTok{"}\NormalTok{, fresh\_var\})}
            \ControlFlowTok{return}\NormalTok{ fresh\_var}

        \ControlFlowTok{else}\NormalTok{:}
            \ControlFlowTok{raise} \PreprocessorTok{NotImplementedError}\NormalTok{(}\StringTok{"Type de noeud non pris en charge."}\NormalTok{)}

\KeywordTok{def}\NormalTok{ print\_cnf(clauses):}
    \CommentTok{"""Affiche de manière lisible la liste des clauses CNF."""}
\NormalTok{    formatted\_clauses }\OperatorTok{=}\NormalTok{ []}
    \ControlFlowTok{for}\NormalTok{ clause }\KeywordTok{in}\NormalTok{ clauses:}
\NormalTok{        sorted\_literals }\OperatorTok{=} \BuiltInTok{sorted}\NormalTok{(}\BuiltInTok{list}\NormalTok{(clause), key}\OperatorTok{=}\KeywordTok{lambda}\NormalTok{ x: x.replace(}\StringTok{"\textasciitilde{}"}\NormalTok{, }\StringTok{""}\NormalTok{))}
\NormalTok{        formatted\_clauses.append(}\StringTok{"("} \OperatorTok{+} \StringTok{" | "}\NormalTok{.join(sorted\_literals) }\OperatorTok{+} \StringTok{")"}\NormalTok{)}
    \BuiltInTok{print}\NormalTok{(}\StringTok{" \& "}\NormalTok{.join(formatted\_clauses))}
\end{lstlisting}

\begin{center}\rule{0.5\linewidth}{0.5pt}\end{center}

\hypertarget{exemple-dutilisation}{%
\subsection{Exemple d'Utilisation}\label{exemple-dutilisation}}

Mettons sous CNF la formule \((P \land Q) \lor R\) :

\begin{lstlisting}
\ControlFlowTok{if} \VariableTok{\_\_name\_\_} \OperatorTok{==} \StringTok{"\_\_main\_\_"}\NormalTok{:}
\NormalTok{    p }\OperatorTok{=}\NormalTok{ Variable(}\StringTok{"P"}\NormalTok{)}
\NormalTok{    q }\OperatorTok{=}\NormalTok{ Variable(}\StringTok{"Q"}\NormalTok{)}
\NormalTok{    r }\OperatorTok{=}\NormalTok{ Variable(}\StringTok{"R"}\NormalTok{)}

    \CommentTok{\# Formule : (P \& Q) | R}
\NormalTok{    formula }\OperatorTok{=}\NormalTok{ Or(And(p, q), r)}
    \BuiltInTok{print}\NormalTok{(}\StringTok{"Formule de départ :"}\NormalTok{, formula)}

\NormalTok{    transformer }\OperatorTok{=}\NormalTok{ TseitinTransformer()}
\NormalTok{    root, clauses }\OperatorTok{=}\NormalTok{ transformer.transform(formula)}

    \BuiltInTok{print}\NormalTok{(}\StringTok{"}\CharTok{\textbackslash{}n}\StringTok{Variable représentative de la racine :"}\NormalTok{, root)}
    \BuiltInTok{print}\NormalTok{(}\StringTok{"}\CharTok{\textbackslash{}n}\StringTok{CNF équisatisfiable générée par Tseitin :"}\NormalTok{)}
\NormalTok{    print\_cnf(clauses)}
\end{lstlisting}

\hypertarget{sortie-console-attendue}{%
\subsubsection{Sortie console attendue
:}\label{sortie-console-attendue}}

\begin{lstlisting}
\NormalTok{Formule de départ : ((P \& Q) | R)}

\NormalTok{Variable représentative de la racine : \_t2}

\NormalTok{CNF équisatisfiable générée par Tseitin :}
\NormalTok{(\textasciitilde{}\_t1 | P) \& (\textasciitilde{}\_t1 | Q) \& (\textasciitilde{}P | \textasciitilde{}Q | \_t1) \& (\textasciitilde{}\_t2 | R | \_t1) \& (\textasciitilde{}\_t1 | \_t2) \& (\textasciitilde{}R | \_t2) \& (\_t2)}
\end{lstlisting}

Ce format est directement exploitable par les solveurs SAT standard.
\n
\hypertarget{tp-4-impluxe9mentation-dun-solveur-sat-dpll}{%
\section{TP 4 : Implémentation d'un solveur SAT
DPLL}\label{tp-4-impluxe9mentation-dun-solveur-sat-dpll}}

\hypertarget{objectif}{%
\subsection{Objectif}\label{objectif}}

L'objectif de ce TP est de concevoir et de coder en Python pur un
\textbf{solveur SAT} complet basé sur l'algorithme historique
\textbf{DPLL} (Davis-Putnam-Logemann-Loveland) incorporant la technique
de la \textbf{propagation unitaire}. Le programme prendra en entrée une
formule sous forme de clauses CNF (comme celle générée au TP 3) et
déterminera si la formule est satisfaisable, en fournissant une
affectation valide des variables le cas échéant.

\begin{center}\rule{0.5\linewidth}{0.5pt}\end{center}

\hypertarget{moduxe9lisation-de-lalgorithme-dpll}{%
\subsection{Modélisation de l'algorithme
DPLL}\label{moduxe9lisation-de-lalgorithme-dpll}}

L'algorithme DPLL est une recherche par séparation et évaluation
(backtracking) dans l'arbre binaire des affectations de variables,
optimisée par la propagation de contraintes forcées.

\hypertarget{les-uxe9tapes-de-la-fonction-ruxe9cursive}{%
\subsubsection{Les étapes de la fonction récursive
:}\label{les-uxe9tapes-de-la-fonction-ruxe9cursive}}

\begin{enumerate}
\def\labelenumi{\arabic{enumi}.}
\tightlist
\item
  \textbf{Propagation unitaire :} Si une clause est réduite à un seul
  littéral \(\{L\}\) (clause unitaire), alors le littéral \(L\) doit
  obligatoirement être évalué à Vrai pour satisfaire cette clause.

  \begin{itemize}
  \tightlist
  \item
    On simplifie les clauses sous l'hypothèse \(L = True\).
  \item
    On répète cette étape jusqu'à ce qu'il n'y ait plus aucune clause
    unitaire.
  \end{itemize}
\item
  \textbf{Conditions d'arrêt :}

  \begin{itemize}
  \tightlist
  \item
    Si l'ensemble des clauses est vide, toutes les contraintes sont
    satisfaites. La formule est \textbf{Satisfaisable}. On renvoie
    \texttt{(True,\ affectation\_courante)}.
  \item
    Si l'ensemble contient une clause vide (un ensemble vide
    \texttt{\{\}}), cela signifie qu'une contradiction a été atteinte.
    La branche actuelle est \textbf{Insatisfaisable}. On renvoie
    \texttt{(False,\ None)}.
  \end{itemize}
\item
  \textbf{Branchement (Heuristique) :}

  \begin{itemize}
  \tightlist
  \item
    On choisit une variable propositionnelle libre \(P\) (qui n'a pas
    encore de valeur affectée).
  \item
    On tente de résoudre récursivement en posant \(P = True\).
  \item
    En cas d'échec (retourne \texttt{False}), on tente de résoudre en
    posant \(P = False\).
  \item
    Si les deux branches échouent, on retourne \texttt{(False,\ None)}.
  \end{itemize}
\end{enumerate}

\begin{center}\rule{0.5\linewidth}{0.5pt}\end{center}

\hypertarget{code-source-python}{%
\subsection{Code Source Python}\label{code-source-python}}

\begin{lstlisting}
\KeywordTok{def}\NormalTok{ parse\_literal(literal):}
    \CommentTok{"""Sépare le signe du nom de la variable. Retourne (nom, est\_negatif)."""}
    \ControlFlowTok{if}\NormalTok{ literal.startswith(}\StringTok{"\textasciitilde{}"}\NormalTok{):}
        \ControlFlowTok{return}\NormalTok{ literal[}\DecValTok{1}\NormalTok{:], }\VariableTok{True}
    \ControlFlowTok{return}\NormalTok{ literal, }\VariableTok{False}

\KeywordTok{def}\NormalTok{ negate\_literal(literal):}
    \CommentTok{"""Retourne le littéral complémentaire."""}
    \ControlFlowTok{if}\NormalTok{ literal.startswith(}\StringTok{"\textasciitilde{}"}\NormalTok{):}
        \ControlFlowTok{return}\NormalTok{ literal[}\DecValTok{1}\NormalTok{:]}
    \ControlFlowTok{return} \SpecialStringTok{f"\textasciitilde{}}\SpecialCharTok{\{}\NormalTok{literal}\SpecialCharTok{\}}\SpecialStringTok{"}

\KeywordTok{def}\NormalTok{ simplify\_cnf(clauses, literal):}
    \CommentTok{"""}
\CommentTok{    Simplifie la CNF en supposant que \textquotesingle{}literal\textquotesingle{} est Vrai.}
\CommentTok{    {-} Supprime toutes les clauses contenant \textquotesingle{}literal\textquotesingle{}.}
\CommentTok{    {-} Supprime la négation de \textquotesingle{}literal\textquotesingle{} de toutes les clauses restantes.}
\CommentTok{    """}
\NormalTok{    neg\_lit }\OperatorTok{=}\NormalTok{ negate\_literal(literal)}
\NormalTok{    new\_clauses }\OperatorTok{=}\NormalTok{ []}

    \ControlFlowTok{for}\NormalTok{ clause }\KeywordTok{in}\NormalTok{ clauses:}
        \ControlFlowTok{if}\NormalTok{ literal }\KeywordTok{in}\NormalTok{ clause:}
            \CommentTok{\# La clause est satisfaite, on la supprime de l\textquotesingle{}ensemble}
            \ControlFlowTok{continue}
        \CommentTok{\# On crée une nouvelle clause en enlevant la négation du littéral}
\NormalTok{        new\_clause }\OperatorTok{=}\NormalTok{ \{lit }\ControlFlowTok{for}\NormalTok{ lit }\KeywordTok{in}\NormalTok{ clause }\ControlFlowTok{if}\NormalTok{ lit }\OperatorTok{!=}\NormalTok{ neg\_lit\}}
\NormalTok{        new\_clauses.append(new\_clause)}

    \ControlFlowTok{return}\NormalTok{ new\_clauses}

\KeywordTok{def}\NormalTok{ dpll(clauses, assignment):}
    \CommentTok{"""}
\CommentTok{    Fonction récursive principale de l\textquotesingle{}algorithme DPLL.}
\CommentTok{    clauses : liste de sets (clauses CNF)}
\CommentTok{    assignment : dictionnaire (variable {-}\textgreater{} bool)}
\CommentTok{    """}
    \CommentTok{\# 1. Propagation unitaire}
\NormalTok{    unit\_clauses }\OperatorTok{=}\NormalTok{ [c }\ControlFlowTok{for}\NormalTok{ c }\KeywordTok{in}\NormalTok{ clauses }\ControlFlowTok{if} \BuiltInTok{len}\NormalTok{(c) }\OperatorTok{==} \DecValTok{1}\NormalTok{]}
    \ControlFlowTok{while}\NormalTok{ unit\_clauses:}
        \CommentTok{\# On extrait le littéral unique de la première clause unitaire}
\NormalTok{        lit }\OperatorTok{=} \BuiltInTok{list}\NormalTok{(unit\_clauses[}\DecValTok{0}\NormalTok{])[}\DecValTok{0}\NormalTok{]}
\NormalTok{        var, is\_neg }\OperatorTok{=}\NormalTok{ parse\_literal(lit)}

        \CommentTok{\# Enregistrement de l\textquotesingle{}affectation forcée}
\NormalTok{        assignment[var] }\OperatorTok{=} \KeywordTok{not}\NormalTok{ is\_neg}

        \CommentTok{\# Simplification de la CNF}
\NormalTok{        clauses }\OperatorTok{=}\NormalTok{ simplify\_cnf(clauses, lit)}

        \CommentTok{\# Cas d\textquotesingle{}arrêt immédiat après simplification}
        \ControlFlowTok{if} \BuiltInTok{len}\NormalTok{(clauses) }\OperatorTok{==} \DecValTok{0}\NormalTok{:}
            \ControlFlowTok{return} \VariableTok{True}\NormalTok{, assignment}
        \ControlFlowTok{if} \BuiltInTok{any}\NormalTok{(}\BuiltInTok{len}\NormalTok{(c) }\OperatorTok{==} \DecValTok{0} \ControlFlowTok{for}\NormalTok{ c }\KeywordTok{in}\NormalTok{ clauses):}
            \ControlFlowTok{return} \VariableTok{False}\NormalTok{, }\VariableTok{None}

\NormalTok{        unit\_clauses }\OperatorTok{=}\NormalTok{ [c }\ControlFlowTok{for}\NormalTok{ c }\KeywordTok{in}\NormalTok{ clauses }\ControlFlowTok{if} \BuiltInTok{len}\NormalTok{(c) }\OperatorTok{==} \DecValTok{1}\NormalTok{]}

    \CommentTok{\# 2. Cas d\textquotesingle{}arrêt de base}
    \ControlFlowTok{if} \BuiltInTok{len}\NormalTok{(clauses) }\OperatorTok{==} \DecValTok{0}\NormalTok{:}
        \ControlFlowTok{return} \VariableTok{True}\NormalTok{, assignment}
    \ControlFlowTok{if} \BuiltInTok{any}\NormalTok{(}\BuiltInTok{len}\NormalTok{(c) }\OperatorTok{==} \DecValTok{0} \ControlFlowTok{for}\NormalTok{ c }\KeywordTok{in}\NormalTok{ clauses):}
        \ControlFlowTok{return} \VariableTok{False}\NormalTok{, }\VariableTok{None}

    \CommentTok{\# 3. Branchement (Choix d\textquotesingle{}une variable non encore assignée)}
    \CommentTok{\# On prend la première variable de la première clause restante}
\NormalTok{    chosen\_lit }\OperatorTok{=} \BuiltInTok{list}\NormalTok{(clauses[}\DecValTok{0}\NormalTok{])[}\DecValTok{0}\NormalTok{]}
\NormalTok{    chosen\_var, \_ }\OperatorTok{=}\NormalTok{ parse\_literal(chosen\_lit)}

    \CommentTok{\# Tentative branche 1 : chosen\_var = True (littéral chosen\_var est vrai)}
\NormalTok{    assignment\_true }\OperatorTok{=}\NormalTok{ assignment.copy()}
\NormalTok{    assignment\_true[chosen\_var] }\OperatorTok{=} \VariableTok{True}
\NormalTok{    clauses\_true }\OperatorTok{=}\NormalTok{ simplify\_cnf(clauses, chosen\_var)}
\NormalTok{    success, res }\OperatorTok{=}\NormalTok{ dpll(clauses\_true, assignment\_true)}
    \ControlFlowTok{if}\NormalTok{ success:}
        \ControlFlowTok{return} \VariableTok{True}\NormalTok{, res}

    \CommentTok{\# Tentative branche 2 : chosen\_var = False (littéral \textasciitilde{}chosen\_var est vrai)}
\NormalTok{    assignment\_false }\OperatorTok{=}\NormalTok{ assignment.copy()}
\NormalTok{    assignment\_false[chosen\_var] }\OperatorTok{=} \VariableTok{False}
\NormalTok{    clauses\_false }\OperatorTok{=}\NormalTok{ simplify\_cnf(clauses, }\SpecialStringTok{f"\textasciitilde{}}\SpecialCharTok{\{}\NormalTok{chosen\_var}\SpecialCharTok{\}}\SpecialStringTok{"}\NormalTok{)}
    \ControlFlowTok{return}\NormalTok{ dpll(clauses\_false, assignment\_false)}

\KeywordTok{def}\NormalTok{ solve\_sat(clauses):}
    \CommentTok{"""Point d\textquotesingle{}entrée utilisateur pour résoudre une CNF."""}
    \CommentTok{\# On s\textquotesingle{}assure d\textquotesingle{}éliminer les doublons et de travailler sur une liste de sets indépendants}
\NormalTok{    cleaned\_clauses }\OperatorTok{=}\NormalTok{ [}\BuiltInTok{set}\NormalTok{(c) }\ControlFlowTok{for}\NormalTok{ c }\KeywordTok{in}\NormalTok{ clauses]}
    \ControlFlowTok{return}\NormalTok{ dpll(cleaned\_clauses, \{\})}
\end{lstlisting}

\begin{center}\rule{0.5\linewidth}{0.5pt}\end{center}

\hypertarget{exemple-de-validation-et-test}{%
\subsection{Exemple de Validation et
Test}\label{exemple-de-validation-et-test}}

Résolvons l'ensemble contradictoire du TP 2 et un cas satisfaisable :

\begin{lstlisting}
\ControlFlowTok{if} \VariableTok{\_\_name\_\_} \OperatorTok{==} \StringTok{"\_\_main\_\_"}\NormalTok{:}
    \CommentTok{\# Test 1 : Base insatisfaisable du TP 2 :}
    \CommentTok{\# (P v Q) \& (\textasciitilde{}P v Q) \& (P v \textasciitilde{}Q) \& (\textasciitilde{}P v \textasciitilde{}Q)}
\NormalTok{    cnf\_contradiction }\OperatorTok{=}\NormalTok{ [}
\NormalTok{        \{}\StringTok{"P"}\NormalTok{, }\StringTok{"Q"}\NormalTok{\},}
\NormalTok{        \{}\StringTok{"\textasciitilde{}P"}\NormalTok{, }\StringTok{"Q"}\NormalTok{\},}
\NormalTok{        \{}\StringTok{"P"}\NormalTok{, }\StringTok{"\textasciitilde{}Q"}\NormalTok{\},}
\NormalTok{        \{}\StringTok{"\textasciitilde{}P"}\NormalTok{, }\StringTok{"\textasciitilde{}Q"}\NormalTok{\}}
\NormalTok{    ]}
    \BuiltInTok{print}\NormalTok{(}\StringTok{"Résolution du Test 1 (Contradiction) :"}\NormalTok{)}
\NormalTok{    sat, model }\OperatorTok{=}\NormalTok{ solve\_sat(cnf\_contradiction)}
    \BuiltInTok{print}\NormalTok{(}\StringTok{"Satisfaisable :"}\NormalTok{, sat)}
    \BuiltInTok{print}\NormalTok{(}\StringTok{"Modèle :"}\NormalTok{, model)}

    \CommentTok{\# Test 2 : Base satisfaisable :}
    \CommentTok{\# (\textasciitilde{}P v Q) \& (\textasciitilde{}Q v R) \& (\textasciitilde{}R v \textasciitilde{}P) \& (P)}
    \CommentTok{\# Doit forcer P=True, ce qui force Q=True (propagation), ce qui force R=True (propagation),}
    \CommentTok{\# ce qui rend la 3ème clause (\textasciitilde{}R v \textasciitilde{}P) fausse car \textasciitilde{}True v \textasciitilde{}True = Faux.}
    \CommentTok{\# On s\textquotesingle{}attend à ce que cette formule soit déclarée insatisfaisable.}
\NormalTok{    cnf\_test2 }\OperatorTok{=}\NormalTok{ [}
\NormalTok{        \{}\StringTok{"\textasciitilde{}P"}\NormalTok{, }\StringTok{"Q"}\NormalTok{\},}
\NormalTok{        \{}\StringTok{"\textasciitilde{}Q"}\NormalTok{, }\StringTok{"R"}\NormalTok{\},}
\NormalTok{        \{}\StringTok{"\textasciitilde{}R"}\NormalTok{, }\StringTok{"\textasciitilde{}P"}\NormalTok{\},}
\NormalTok{        \{}\StringTok{"P"}\NormalTok{\}}
\NormalTok{    ]}
    \BuiltInTok{print}\NormalTok{(}\StringTok{"}\CharTok{\textbackslash{}n}\StringTok{Résolution du Test 2 (Propagation forcée menant à un échec) :"}\NormalTok{)}
\NormalTok{    sat, model }\OperatorTok{=}\NormalTok{ solve\_sat(cnf\_test2)}
    \BuiltInTok{print}\NormalTok{(}\StringTok{"Satisfaisable :"}\NormalTok{, sat)}
    \BuiltInTok{print}\NormalTok{(}\StringTok{"Modèle :"}\NormalTok{, model)}

    \CommentTok{\# Test 3 : Base satisfaisable :}
    \CommentTok{\# (\textasciitilde{}P v Q) \& (\textasciitilde{}Q v R) \& (P)}
    \CommentTok{\# Doit forcer P=True, Q=True, R=True}
\NormalTok{    cnf\_satisfaisable }\OperatorTok{=}\NormalTok{ [}
\NormalTok{        \{}\StringTok{"\textasciitilde{}P"}\NormalTok{, }\StringTok{"Q"}\NormalTok{\},}
\NormalTok{        \{}\StringTok{"\textasciitilde{}Q"}\NormalTok{, }\StringTok{"R"}\NormalTok{\},}
\NormalTok{        \{}\StringTok{"P"}\NormalTok{\}}
\NormalTok{    ]}
    \BuiltInTok{print}\NormalTok{(}\StringTok{"}\CharTok{\textbackslash{}n}\StringTok{Résolution du Test 3 (Satisfaisable avec propagation) :"}\NormalTok{)}
\NormalTok{    sat, model }\OperatorTok{=}\NormalTok{ solve\_sat(cnf\_satisfaisable)}
    \BuiltInTok{print}\NormalTok{(}\StringTok{"Satisfaisable :"}\NormalTok{, sat)}
    \BuiltInTok{print}\NormalTok{(}\StringTok{"Modèle :"}\NormalTok{, model)}
\end{lstlisting}

\hypertarget{sortie-attendue}{%
\subsubsection{Sortie attendue :}\label{sortie-attendue}}

\begin{lstlisting}
\NormalTok{Résolution du Test 1 (Contradiction) :}
\NormalTok{Satisfaisable : False}
\NormalTok{Modèle : None}

\NormalTok{Résolution du Test 2 (Propagation forcée menant à un échec) :}
\NormalTok{Satisfaisable : False}
\NormalTok{Modèle : None}

\NormalTok{Résolution du Test 3 (Satisfaisable avec propagation) :}
\NormalTok{Satisfaisable : True}
\NormalTok{Modèle : \{\textquotesingle{}P\textquotesingle{}: True, \textquotesingle{}Q\textquotesingle{}: True, \textquotesingle{}R\textquotesingle{}: True\}}
\end{lstlisting}

Our SAT-solver is fully operational!
\n
\hypertarget{tp-5-ruxe9solution-du-probluxe8me-des-n-reines-par-ruxe9duction-sat}{%
\section{TP 5 : Résolution du problème des N-Reines par réduction
SAT}\label{tp-5-ruxe9solution-du-probluxe8me-des-n-reines-par-ruxe9duction-sat}}

\hypertarget{objectif}{%
\subsection{Objectif}\label{objectif}}

L'objectif de ce TP est d'utiliser le solveur SAT DPLL développé au TP 4
pour résoudre un problème combinatoire classique d'Intelligence
Artificielle : \textbf{le problème des N-Reines} (placer \(N\) reines
sur un échiquier de taille \(N \times N\) sans qu'elles ne se menacent
mutuellement). Nous écrirons un script Python pour générer
automatiquement les contraintes du problème sous forme de clauses CNF,
les résoudrons avec notre solveur et afficherons la grille solution.

\begin{center}\rule{0.5\linewidth}{0.5pt}\end{center}

\hypertarget{moduxe9lisation-du-probluxe8me-des-n-reines-en-sat}{%
\subsection{Modélisation du problème des N-Reines en
SAT}\label{moduxe9lisation-du-probluxe8me-des-n-reines-en-sat}}

Pour un échiquier de taille \(N \times N\), nous définissons \(N^2\)
variables propositionnelles notées \texttt{Q\_i\_j} pour
\(0 \le i, j < N\). - La variable \texttt{Q\_i\_j} est évaluée à Vrai si
et seulement si une reine est placée à la ligne \(i\) et à la colonne
\(j\).

\hypertarget{les-contraintes-du-probluxe8me}{%
\subsubsection{Les Contraintes du Problème
:}\label{les-contraintes-du-probluxe8me}}

\begin{enumerate}
\def\labelenumi{\arabic{enumi}.}
\tightlist
\item
  \textbf{Au moins une reine par ligne :} Pour chaque ligne \(i\) :
  \[\bigvee_{j=0}^{N-1} Q_{i, j}\]
\item
  \textbf{Au plus une reine par ligne :} Pour chaque ligne \(i\), et
  chaque paire de colonnes distinctes \(j_1 < j_2\) :
  \[\neg Q_{i, j_1} \lor \neg Q_{i, j_2}\]
\item
  \textbf{Au plus une reine par colonne :} Pour chaque colonne \(j\), et
  chaque paire de lignes distinctes \(i_1 < i_2\) :
  \[\neg Q_{i_1, j} \lor \neg Q_{i_2, j}\]
\item
  \textbf{Au plus une reine par diagonale :} Deux reines sur une même
  diagonale (descendante ou montante) se menacent si la valeur absolue
  de la différence de leurs lignes est égale à celle de leurs colonnes.
  Pour toute paire de cases \((i_1, j_1)\) et \((i_2, j_2)\) telles que
  \(i_1 < i_2\) et \(|i_1 - i_2| = |j_1 - j_2|\) :
  \[\neg Q_{i_1, j_1} \lor \neg Q_{i_2, j_2}\]
\end{enumerate}

\begin{center}\rule{0.5\linewidth}{0.5pt}\end{center}

\hypertarget{code-source-python}{%
\subsection{Code Source Python}\label{code-source-python}}

\begin{lstlisting}
\CommentTok{\# Import du solveur du TP 4}
\ImportTok{from}\NormalTok{ tp\_04\_dpll }\ImportTok{import}\NormalTok{ solve\_sat}

\KeywordTok{def}\NormalTok{ get\_var\_name(i, j):}
    \CommentTok{"""Retourne le nom de la variable propositionnelle pour la case (i, j)."""}
    \ControlFlowTok{return} \SpecialStringTok{f"Q\_}\SpecialCharTok{\{}\NormalTok{i}\SpecialCharTok{\}}\SpecialStringTok{\_}\SpecialCharTok{\{}\NormalTok{j}\SpecialCharTok{\}}\SpecialStringTok{"}

\KeywordTok{def}\NormalTok{ generate\_n\_queens\_cnf(n):}
    \CommentTok{"""Génère la liste des clauses CNF pour le problème des N{-}Reines."""}
\NormalTok{    clauses }\OperatorTok{=}\NormalTok{ []}

    \CommentTok{\# 1. Au moins une reine par ligne}
    \ControlFlowTok{for}\NormalTok{ i }\KeywordTok{in} \BuiltInTok{range}\NormalTok{(n):}
\NormalTok{        clause }\OperatorTok{=}\NormalTok{ \{get\_var\_name(i, j) }\ControlFlowTok{for}\NormalTok{ j }\KeywordTok{in} \BuiltInTok{range}\NormalTok{(n)\}}
\NormalTok{        clauses.append(clause)}

    \CommentTok{\# 2. Au plus une reine par ligne}
    \ControlFlowTok{for}\NormalTok{ i }\KeywordTok{in} \BuiltInTok{range}\NormalTok{(n):}
        \ControlFlowTok{for}\NormalTok{ j1 }\KeywordTok{in} \BuiltInTok{range}\NormalTok{(n):}
            \ControlFlowTok{for}\NormalTok{ j2 }\KeywordTok{in} \BuiltInTok{range}\NormalTok{(j1 }\OperatorTok{+} \DecValTok{1}\NormalTok{, n):}
\NormalTok{                clauses.append(\{}\SpecialStringTok{f"\textasciitilde{}}\SpecialCharTok{\{}\NormalTok{get\_var\_name(i, j1)}\SpecialCharTok{\}}\SpecialStringTok{"}\NormalTok{, }\SpecialStringTok{f"\textasciitilde{}}\SpecialCharTok{\{}\NormalTok{get\_var\_name(i, j2)}\SpecialCharTok{\}}\SpecialStringTok{"}\NormalTok{\})}

    \CommentTok{\# 3. Au plus une reine par colonne}
    \ControlFlowTok{for}\NormalTok{ j }\KeywordTok{in} \BuiltInTok{range}\NormalTok{(n):}
        \ControlFlowTok{for}\NormalTok{ i1 }\KeywordTok{in} \BuiltInTok{range}\NormalTok{(n):}
            \ControlFlowTok{for}\NormalTok{ i2 }\KeywordTok{in} \BuiltInTok{range}\NormalTok{(i1 }\OperatorTok{+} \DecValTok{1}\NormalTok{, n):}
\NormalTok{                clauses.append(\{}\SpecialStringTok{f"\textasciitilde{}}\SpecialCharTok{\{}\NormalTok{get\_var\_name(i1, j)}\SpecialCharTok{\}}\SpecialStringTok{"}\NormalTok{, }\SpecialStringTok{f"\textasciitilde{}}\SpecialCharTok{\{}\NormalTok{get\_var\_name(i2, j)}\SpecialCharTok{\}}\SpecialStringTok{"}\NormalTok{\})}

    \CommentTok{\# 4. Au plus une reine par diagonale}
    \ControlFlowTok{for}\NormalTok{ i1 }\KeywordTok{in} \BuiltInTok{range}\NormalTok{(n):}
        \ControlFlowTok{for}\NormalTok{ j1 }\KeywordTok{in} \BuiltInTok{range}\NormalTok{(n):}
            \ControlFlowTok{for}\NormalTok{ i2 }\KeywordTok{in} \BuiltInTok{range}\NormalTok{(i1 }\OperatorTok{+} \DecValTok{1}\NormalTok{, n):}
                \ControlFlowTok{for}\NormalTok{ j2 }\KeywordTok{in} \BuiltInTok{range}\NormalTok{(n):}
                    \ControlFlowTok{if} \BuiltInTok{abs}\NormalTok{(i1 }\OperatorTok{{-}}\NormalTok{ i2) }\OperatorTok{==} \BuiltInTok{abs}\NormalTok{(j1 }\OperatorTok{{-}}\NormalTok{ j2):}
\NormalTok{                        clauses.append(\{}\SpecialStringTok{f"\textasciitilde{}}\SpecialCharTok{\{}\NormalTok{get\_var\_name(i1, j1)}\SpecialCharTok{\}}\SpecialStringTok{"}\NormalTok{, }\SpecialStringTok{f"\textasciitilde{}}\SpecialCharTok{\{}\NormalTok{get\_var\_name(i2, j2)}\SpecialCharTok{\}}\SpecialStringTok{"}\NormalTok{\})}

    \ControlFlowTok{return}\NormalTok{ clauses}

\KeywordTok{def}\NormalTok{ print\_board(n, model):}
    \CommentTok{"""Affiche l\textquotesingle{}échiquier avec les reines placées."""}
\NormalTok{    border }\OperatorTok{=} \StringTok{"+"} \OperatorTok{+} \StringTok{"{-}{-}{-}+"} \OperatorTok{*}\NormalTok{ n}
    \BuiltInTok{print}\NormalTok{(border)}
    \ControlFlowTok{for}\NormalTok{ i }\KeywordTok{in} \BuiltInTok{range}\NormalTok{(n):}
\NormalTok{        row }\OperatorTok{=} \StringTok{"|"}
        \ControlFlowTok{for}\NormalTok{ j }\KeywordTok{in} \BuiltInTok{range}\NormalTok{(n):}
\NormalTok{            var }\OperatorTok{=}\NormalTok{ get\_var\_name(i, j)}
            \CommentTok{\# Si la variable vaut True dans le modèle, on place une reine \textquotesingle{}R\textquotesingle{}}
            \ControlFlowTok{if}\NormalTok{ model.get(var, }\VariableTok{False}\NormalTok{):}
\NormalTok{                row }\OperatorTok{+=} \StringTok{" R |"}
            \ControlFlowTok{else}\NormalTok{:}
\NormalTok{                row }\OperatorTok{+=} \StringTok{" . |"}
        \BuiltInTok{print}\NormalTok{(row)}
        \BuiltInTok{print}\NormalTok{(border)}

\KeywordTok{def}\NormalTok{ solve\_n\_queens(n):}
    \CommentTok{"""Génère les contraintes, lance le solveur et affiche le résultat."""}
    \BuiltInTok{print}\NormalTok{(}\SpecialStringTok{f"{-}{-}{-} Résolution du problème des }\SpecialCharTok{\{}\NormalTok{n}\SpecialCharTok{\}}\SpecialStringTok{{-}Reines {-}{-}{-}"}\NormalTok{)}

    \CommentTok{\# Génération des contraintes}
\NormalTok{    clauses }\OperatorTok{=}\NormalTok{ generate\_n\_queens\_cnf(n)}
    \BuiltInTok{print}\NormalTok{(}\SpecialStringTok{f"Nombre de clauses CNF générées : }\SpecialCharTok{\{}\BuiltInTok{len}\NormalTok{(clauses)}\SpecialCharTok{\}}\SpecialStringTok{"}\NormalTok{)}

    \CommentTok{\# Résolution SAT}
\NormalTok{    success, model }\OperatorTok{=}\NormalTok{ solve\_sat(clauses)}

    \ControlFlowTok{if}\NormalTok{ success:}
        \BuiltInTok{print}\NormalTok{(}\StringTok{"}\CharTok{\textbackslash{}n}\StringTok{Solution trouvée :"}\NormalTok{)}
\NormalTok{        print\_board(n, model)}
    \ControlFlowTok{else}\NormalTok{:}
        \BuiltInTok{print}\NormalTok{(}\StringTok{"}\CharTok{\textbackslash{}n}\StringTok{Aucune solution n\textquotesingle{}existe pour cette taille d\textquotesingle{}échiquier."}\NormalTok{)}
\end{lstlisting}

\begin{center}\rule{0.5\linewidth}{0.5pt}\end{center}

\hypertarget{exemple-de-validation-et-test-n-4}{%
\subsection{Exemple de Validation et Test (N =
4)}\label{exemple-de-validation-et-test-n-4}}

Placer 4 reines sur un échiquier de 4x4 :

\begin{lstlisting}
\ControlFlowTok{if} \VariableTok{\_\_name\_\_} \OperatorTok{==} \StringTok{"\_\_main\_\_"}\NormalTok{:}
    \CommentTok{\# Résolution pour N = 4}
\NormalTok{    solve\_n\_queens(}\DecValTok{4}\NormalTok{)}

    \CommentTok{\# Résolution pour N = 3 (doit être déclarée impossible)}
    \BuiltInTok{print}\NormalTok{(}\StringTok{"}\CharTok{\textbackslash{}n}\StringTok{"}\NormalTok{)}
\NormalTok{    solve\_n\_queens(}\DecValTok{3}\NormalTok{)}
\end{lstlisting}

\hypertarget{rendu-console-attendu}{%
\subsubsection{Rendu console attendu :}\label{rendu-console-attendu}}

\begin{lstlisting}
\NormalTok{{-}{-}{-} Résolution du problème des 4{-}Reines {-}{-}{-}}
\NormalTok{Nombre de clauses CNF générées : 84}

\NormalTok{Solution trouvée :}
\NormalTok{+{-}{-}{-}+{-}{-}{-}+{-}{-}{-}+{-}{-}{-}+}
\NormalTok{| . | R | . | . |}
\NormalTok{+{-}{-}{-}+{-}{-}{-}+{-}{-}{-}+{-}{-}{-}+}
\NormalTok{| . | . | . | R |}
\NormalTok{+{-}{-}{-}+{-}{-}{-}+{-}{-}{-}+{-}{-}{-}+}
\NormalTok{| R | . | . | . |}
\NormalTok{+{-}{-}{-}+{-}{-}{-}+{-}{-}{-}+{-}{-}{-}+}
\NormalTok{| . | . | R | . |}
\NormalTok{+{-}{-}{-}+{-}{-}{-}+{-}{-}{-}+{-}{-}{-}+}


\NormalTok{{-}{-}{-} Résolution du problème des 3{-}Reines {-}{-}{-}}
\NormalTok{Nombre de clauses CNF générées : 26}

\NormalTok{Aucune solution n\textquotesingle{}existe pour cette taille d\textquotesingle{}échiquier.}
\end{lstlisting}

Le placement des reines est correct, les diagonales, lignes et colonnes
sont parfaitement préservées !
\n


\end{document}